\documentclass[11pt,letterpaper]{article}
\usepackage[T1]{fontenc}
\usepackage[utf8]{inputenc}
\usepackage{lmodern,microtype}
\usepackage[margin=1.01in,top=.9in,bottom=.91in,headheight=15pt]{geometry}
\usepackage{amsmath,amssymb,amsthm,mathtools,booktabs,longtable,tabularx,array,enumitem,fancyhdr,xcolor,xurl}
\usepackage[colorlinks=true,linkcolor=teal!55!black,urlcolor=teal!65!black,citecolor=teal!55!black]{hyperref}
\usepackage{bookmark}
\definecolor{ink}{HTML}{17374A}
\definecolor{accent}{HTML}{116661}
\hypersetup{pdftitle={First-order logic in intuitionistic primitive recursive arithmetic},pdfsubject={Proof transformations, finite models, decision procedures and limits},pdfauthor={Research report prepared for Vladimir Reshetnikov}}
\pagestyle{fancy}\fancyhf{}\fancyhead[L]{\small\sffamily\color{ink}First-order logic in iPRA}\fancyhead[R]{\small\sffamily\color{ink}Proofs and finite models}\fancyfoot[C]{\small\thepage}\renewcommand{\headrulewidth}{.25pt}
\newtheoremstyle{packet}{8pt}{4pt}{\itshape}{}{\bfseries\color{ink}}{.}{.5em}{}
\theoremstyle{packet}
\newtheorem{theorem}{Theorem}[section]
\newtheorem{lemma}[theorem]{Lemma}
\newtheorem{corollary}[theorem]{Corollary}
\newtheorem{proposition}[theorem]{Proposition}
\theoremstyle{definition}\newtheorem{definition}[theorem]{Definition}\newtheorem{example}[theorem]{Example}
\setlength{\emergencystretch}{3em}
\setlist{topsep=3pt,itemsep=2pt}
\setcounter{tocdepth}{1}
\hypersetup{bookmarksdepth=2}
\makeatletter
\renewcommand*\l@section[2]{\addpenalty\@secpenalty\vskip .21em\@dottedtocline{1}{0em}{2.1em}{\bfseries #1}{#2}}
\makeatother
\begin{document}
\thispagestyle{plain}
{\noindent\sffamily\small\color{accent}MATHEMATICAL LOGIC\quad /\quad CONSTRUCTIVE METATHEORY\par}
\vspace{1.15em}
{\raggedright\noindent\color{ink}\huge\bfseries First-order logic in intuitionistic\\primitive recursive arithmetic\par}
\vspace{.6em}
{\noindent\large Proof transformations, finite models and limits of decision\par}
\vspace{1.2em}
\noindent A first-order formula, a derivation and a finite interpretation can each be written as a natural number. This lets weak arithmetic prove theorems about whole families of formulas, including classical formulas. The essential question is what data each theorem produces. A negative translation produces a new proof. A failed finite back-and-forth test produces a distinguishing formula. A successful Skolemization of a finite model produces actual tables of witnesses. None of these constructions needs an oracle for truth.
\medskip

\noindent Two examples locate the boundary. With $p$ unary predicates and equality, the truth of a sentence of quantifier rank $q$ on finite structures is determined by $2^p$ counts truncated at $q$. With a single binary relation one can encode a computation, and even a procedure that only separates provable sentences from those with finite countermodels is impossible. Both arguments give explicit algorithms on the side where an algorithm exists.
\medskip

\noindent The manuscript gives human proofs in iPRA, a Lean formalization plan tied to ProveIt's ordinary iPRA calculus, and a smaller kernel-checked Lean development. The implementation and verification statuses are recorded in Section~\ref{sec:lean}.
\tableofcontents\clearpage
\section{What iPRA must prove}\label{sec:arith}
Our base theory is intuitionistic first-order arithmetic with $0,S$, injectivity of $S$, $Sx\ne0$, function symbols for primitive-recursive definitions with their defining equations, and induction on quantifier-free formulas with parameters. A displayed theorem about codes abbreviates its universal closure with well-formedness hypotheses. A new algorithm stands for a specific finite primitive-recursive (PR) definition, or for a schema in a named PR parameter. An arbitrary algorithm given by its program index is run only for a supplied number of steps. This convention avoids a universal evaluator for all PR definitions.

\begin{lemma}[A coding and induction toolkit]\label{lem:toolkit}
Finite arrays, bounded search with a default, finite conjunction, and bounded verification of histories have PR definitions whose elementary correctness laws are provable in iPRA. Primitive recursion on an externally supplied height, with a fixed PR bound per stage, is permitted even when the resulting numbers are very large.
\end{lemma}
\begin{proof}
One concrete pairing is $P(a,b)=2^a(2b+1)$. Reserve $0$ for the empty list and cons by $P(a,b)$ with $b$ the tail. Each component is strictly smaller than the positive pair. To decode $u>0$, strip factors of two at most $u$ times while the remainder is even, counting divisions. This terminates with an odd positive number: every active division strictly decreases the remainder, so a still-active run of $u$ steps is impossible by the invariant that the remainder after $j$ active steps is at most $u-j$. Halving, parity and comparison are PR by counting. Induction gives $2^a(2b+1)$ back from the projections, and stripping a constructed pair gives exactly $a,b$. Traversing a coded list now takes at most its code in steps. Append, reverse, lookup with default, and bounded replacement are iterations. One may instead materialize tables in fixed radix $B\ge2$, using $\sum_{i<L}v_i B^i$ when all entries are less than $B$.
 
For a bit test $h(i)$ fold AND or OR over $i<b$. A search stores the first success, with $b$ as sentinel. At stage $j$ maintain the bit saying either the stored index is below $j$, satisfies the test and no earlier one does, or all the first $j$ tests fail and the sentinel remains. Prefix AND gives the last two bounded conditions. The step checks the new bit. Pointwise correctness of the fold can also be proved by propagating an arbitrary fixed counterindex through prefixes or recovering it by the same search. Thus the bounded universal or existential may be freely expanded in the outer logic. The induction formula itself is an equation between numbers. Bounded finite choice is obtained by running the least-witness search independently in each row.

Given a deterministic loop of $l$ iterations with PR step, code its successive states by primitive recursion. A PR bound $F$ on a single update can be iterated $l$ times to bound the largest state (replace $F$ by a monotone majorant using finite maxima if necessary). For proofs by recursion over trees use a finite list of local obligations. Compute child arguments downwards, then verify the local equalities or implications upwards. Store a bit for the conjunction of the verified rows up to the current index. Starting with the empty conjunction, a new row follows from its already verified children by its rule. The logical proof that a bit equals one can use parameters freely. For branching under a quantifier over a domain of size $d$, allocate all $d$ possible calls. For a path of depth $h$ with at most $b$ successors the list has at most $(b+1)^{h+1}$ entries. For a varying bound take the maximum over the stage bounds just described. This method licenses the structural inductions used below when a substitution, assignment or strategy changes between calls. It does not turn a structural induction with an arbitrary unbounded invariant into quantifier-free induction.
\end{proof}

\paragraph{Three readings of a formula.}
$\operatorname{Check}_I(p,\Gamma,A)=1$ says a finite tree uses the intuitionistic object rules. A separate tag permits classical double-negation elimination. $V(M,s,A)=1$ is a numerical truth test in a nonempty, explicitly enumerated, two-valued structure. Forcing in a finite Kripke structure gives a different test. The ambient iPRA proves equations and implications about all these tests. Thus, for example, a truth-table proof of excluded middle for $V$ gives no rule for eliminating double negation in iPRA. Statements about an arbitrary set-sized structure are identified explicitly as external readings. All numerical model statements quantify over codes of valid tables.

\paragraph{How to locate the source claims.}
Formation, variables, interpretations, deductive systems, equality, decidable fragments, expressive limitations and variations are headings in the starting article. Its linked pages on completeness, compactness, Ehrenfeucht--Fraïssé games, prenex and Skolem normal forms, relational algebra and extensions by definitions motivate the corresponding theorems. The source guide at the end separates these connections from additional lemmas needed in the proofs.

\section{Syntax with a computable account of binding}\label{sec:binding}

Fix a finite signature, given by two lists of arities for functions and
predicates. Constants have function arity zero. The arguments also apply
schema by schema to a fixed signature given by primitive-recursive (PR)
symbol tests and arity functions. An expression at arity $n$ may refer to
the $n$ individuals $v_0,\ldots,v_{n-1}$. Terms are variables and
applications $f(\vec t)$ of the stipulated arity. Formulas are predicates
$R(\vec t)$, identities $s=t$, $\bot$, $A\wedge B$, $A\vee B$, $A\to B$,
$\forall B$ and $\exists B$. Under either quantifier the body has arity
$n+1$, with its newly bound slot at index zero. Write
$\neg A=A\to\bot$ and $A\leftrightarrow B=(A\to B)\wedge(B\to A)$.
An expression remains well scoped if an unused slot is appended at
the high end. A sentence has no free occurrences.

\begin{theorem}[Parsing, scope and alpha-conversion]\label{thm:alpha}
Well-formedness, the parse tree, the free-variable set, the binding
occurrence of each bound variable and alpha-equivalence of named
expressions all have PR algorithms. Parsing succeeds with a unique tree.
\end{theorem}
\begin{proof}
For a concrete serialization use prefix tokens carrying the constructor
and any index or symbol. Its arity fixes the number and sorts of its
children. Read the tokens backwards, pushing leaves and at each internal
node popping the prescribed number of children, checking sorts in their
specified order and pushing the resulting tree. Reject underflow,
a bad symbol or any final stack other than one element of the requested
sort. At each suffix the stack is its left-to-right list of complete
subtrees whenever that suffix can arise in a parse. Reading a parent is
then forced. Induction on the number of scanned tokens proves uniqueness
and proves that serializing any tree is inverted by this scan. To check
indices, traverse the resulting tree with the current binder depth $d$.
An index $j<d$ picks the binder $j$ positions back from the innermost.
Otherwise it denotes free slot $j-d$ and must satisfy $j-d<n$.
Mark those slots in an $n$-bit mask. This computes the exact support,
so closedness is decidable even when the input was supplied at positive
arity.

For named input keep a stack of active binder names, nearest first.
Look up each occurrence on that stack by bounded search. A hit gives its
position, and a miss gives the free name's index in a fixed table shared
by both expressions being compared. A common table is obtained by sorting
and removing duplicates from the free names in the two inputs. Renaming a
binder to a fresh name and changing exactly its bound occurrences preserves
these lookups. Conversely, assign to each binder in preorder the name
$M+1+j$ for its preorder number $j$, where $M$ exceeds all names in both
inputs. Both expressions can be renamed to this representative one binder
at a time without capture. Equal indexed trees now print identically.
The scans, stacks and bounded searches use at most the total token count
in steps per traversal or lookup. The fresh names are bounded by $M+1$
plus the larger token count.
\end{proof}

A substitution $\sigma:n\to m$ is a list of $n$ terms at arity $m$.
Let $E[\sigma]$ replace each free slot by its image. On terms the
clauses are $v_i[\sigma]=\sigma_i$ and
$f(\vec t)[\sigma]=f(\vec t[\sigma])$. For atomic formulas and
connectives apply the operation to the children. Put
$\uparrow t=t[v_1,\ldots,v_m]$ for a term at arity $m$, and use
the same simultaneous shift on formulas. At a quantifier use the lift
\[
 \sigma^+=(v_0,\uparrow\sigma_0,\ldots,\uparrow\sigma_{n-1}),
 \qquad (QB)[\sigma]=Q(B[\sigma^+]).
\]
Equivalently, at depth $d$ copy $v_j$ when $j<d$, and replace $v_{d+i}$
by $\sigma_i$ with every index increased by $d$. This last algorithm
defines shift on formulas directly. Write
$B\langle t\rangle=B[t,v_0,\ldots,v_{n-1}]$ for opening an
$(n+1)$-ary body at arity $n$.

\begin{theorem}[Simultaneous substitution]\label{thm:subst}
Substitution preserves scope, acts identically for the variable list
$\mathrm{id}_n$, and obeys
\[
  (E[\sigma])[\tau]=E[\sigma;\tau],\qquad
  (\sigma;\tau)_i=\sigma_i[\tau].
\]
Its exact free support is
$\operatorname{FV}(E[\sigma])
 =\bigcup_{i\in\operatorname{FV}(E)}\operatorname{FV}(\sigma_i)$.
Substitutions agreeing on the free support give identical results.
In particular
$(B\langle t\rangle)[\sigma]=B[\sigma^+]\langle t[\sigma]\rangle$.
\end{theorem}
\begin{proof}
At depth $d$ the copied bound indices are below $d$ and every inserted
index lies in $[d,d+m)$. A former free occurrence contributes exactly
the free variables of its replacement. Taking the union over the
occurrences gives the support formula, and comparison of the actual
insertions proves agreement. Replacing by the identity copies every
index. For composition consider a free occurrence $v_{d+i}$. After
the first operation its replacement is $\sigma_i$ shifted by $d$.
Every variable in this inserted term is therefore free relative to
the surrounding $d$ binders. Applying the second operation replaces
each original variable $v_j$ of $\sigma_i$ by $\tau_j$ shifted
by $d$. The result is precisely $\sigma_i[\tau]$ shifted by $d$.
Original bound occurrences and constructor nodes are copied on both
sides. This proves the equality of output trees, including when
$\sigma_i$ has no variables. For opening, the two composites send
$v_0$ to $t[\sigma]$ and $v_{i+1}$ to $\sigma_i$. Apply composition.

For explicit PR bounds, if $N$ counts input nodes and $L$ is the
largest image size (take $L=1$ for an empty list), the output has at
most $N\max(1,L)$ nodes. The depth is at most $N$, and a shifted
variable index is less than $m+N$ whenever one is inserted. A traversal
with a stack visits each original node and copies each selected image
within these bounds. The preceding arguments can be verified by a
bounded conjunction over occurrence addresses: trace each output block
to its original node, then use the computed input path and depth at that
node. This avoids induction with a universally quantified substitution
as induction formula.
\end{proof}

\section{Effective deduction and admissible proof operations}\label{sec:deduction}

Write $\Gamma\vdash_n A$ for a certificate in arity $n$ with an
ordered finite list $\Gamma$ of assumptions. Write $\Gamma,A$ for
appending an assumption, and $\uparrow\Gamma$ for shifting every
formula. A node records its sequent, a rule, its ordered child proofs
and the term or body used by its rule. All formulas and terms are
scope-checked. The two calculi below share the following rules.
\begin{itemize}
\item A leaf cites an assumption by position. Ex falso concludes
any $A$ from $\bot$. Conjunction is formed from both proofs and projected
to either component. Disjunction is introduced from either disjunct.
To eliminate $A\vee B$, prove $C$ under $\Gamma,A$ and under
$\Gamma,B$. Implication introduction discharges the last assumption
in $\Gamma,A\vdash B$ and elimination is modus ponens in a common
context.
\item Universal introduction has premise
$\uparrow\Gamma\vdash_{n+1}B$ and conclusion
$\Gamma\vdash_n\forall B$. Its elimination gives $B\langle t\rangle$
from $\forall B$. Existential introduction reverses this last pattern
with $\exists$ in the conclusion. Existential elimination has premises
$\Gamma\vdash_n\exists B$ and
$\uparrow\Gamma,B\vdash_{n+1}\uparrow C$, and conclusion
$\Gamma\vdash_n C$.
\item Reflexivity proves $t=t$. Equality elimination takes proofs
of $s=t$ and $B\langle s\rangle$ to a proof of $B\langle t\rangle$.
The unary formula context $B$ can also use the old $n$ slots.
\item A fresh, unused parameter may be dropped:
from $\uparrow\Gamma\vdash_{n+1}\uparrow A$ infer
$\Gamma\vdash_n A$. This implements the ordinary convention of
nonempty domains even in a signature without a closed term. For example,
in one parameter, reflexivity followed by existential introduction proves
the shift of $\exists x\,(x=x)$. The rule closes this proof.
\end{itemize}
Call this calculus $I$. Its extension $C$ adds one rule, from
$\neg\neg A$ infer $A$. The tags $I,C$ here describe object proofs.
All the algorithms and arguments about their codes use iPRA.

\begin{theorem}[Proof recognition and finite search]\label{thm:proofcheck}
For either calculus, the relation $\operatorname{Check}(p,n,\Gamma,A)$
is PR. The increasing bounded search
$S(b)=\bigvee_{p\le b}\operatorname{Check}(p,n,\Gamma,A)$ is PR,
and a positive result supplies a least certificate $p\le b$.
Consequently theoremhood in either presented calculus is recursively
enumerable.
\end{theorem}
\begin{proof}
Parse a proof code as a tree. Reject malformed input. At each node compare
its child sequents to the required instances, computing shifts, openings
and positional lookups. Mark it accepted exactly when these tests and
all child marks succeed. A bottom-up scan proves by induction over the
scan position that the accepted nodes are exactly the locally justified
subtrees with accepted children. The root must also match the four inputs.
Operations use a traversal bounded by input size and the PR substitution
algorithm. Bounded disjunction and least successful index are finite
folds. Enumerate proof codes by increasing number and output their
accepted sequents. Each actual proof occurs at its own code number.
\end{proof}

\begin{theorem}[Instance, context substitution, cut and deduction]\label{thm:graft}
Uniform PR operations on certificates implement all the following.
\begin{enumerate}
\item From $\Gamma\vdash_n A$ and $\sigma:n\to m$, obtain
$\Gamma[\sigma]\vdash_m A[\sigma]$.
\item Given $p:\Gamma\vdash_n A$ and, for every position $i$ of
$\Gamma$, a proof $q_i:\Delta\vdash_n\Gamma_i$, obtain
$\Delta\vdash_n A$. In particular weakening, exchange, contraction
and cut are admissible.
\item Certificates for $\Gamma,A\vdash B$ and for
$\Gamma\vdash A\to B$ can be converted in both directions.
A derivation uses only a computable sublist of its assumptions.
\end{enumerate}
\end{theorem}
\begin{proof}
For the first operation relabel every sequent and the rule's term and
formula annotations by $\sigma$. At an individual-variable binder use
$\sigma^+$ in the premise at larger arity. In initial contexts and
shifted conclusions we require
$(\uparrow E)[\sigma^+]=\uparrow(E[\sigma])$. The two composite
variable maps coincide by Theorem~\ref{thm:subst}. In instantiation
and in equality elimination we require that opening commute with
substitution, the last identity of that theorem. Thus the same rule
tag remains legal everywhere, including at parameter erasure.
Propositional binders merely retain their added assumption.

For the second operation, first implement change of context by a map
from each assumption position to a position carrying the same formula.
Change hypothesis leaves by the map. Under a discharged assumption
extend the map by matching the new last position to the new last
position. Under an individual binder the old position map also maps the
shifted formulas. For existential elimination extend that shifted map
by the new last position. All other steps copy their original rule.
This gives weakening and all positional structural rules directly.

Now graft $q_i$ at each original assumption leaf. On entering a
propositional discharge with new assumption $D$, weaken every current
$q_i$ to $\Delta,D$ and append the assumption leaf for $D$ to the
substitution list. At universal introduction or parameter erasure,
shift the current proofs by part 1. At the second child of existential
elimination do both operations, first shift and then append the proof
of its new local assumption. All premise sequents are thereby relabelled
as required. Cut is the case of one nontrivial graft and identity
assumption proofs for the other positions. Discharge gives the forward
deduction conversion in a single step. For the reverse direction weaken
the proof of $A\to B$ to $\Gamma,A$, append an assumption leaf for
$A$ and use modus ponens.

To extract the finite support, trace the cited hypothesis of each
assumption leaf back through the position maps to the root or to its
local discharge. Mark precisely the cited root positions. Delete the
other root positions throughout and reindex surviving citations and local
discharges. This yields a proof under the marked sublist. In particular,
a proof from a possibly infinite axiom collection, presented with a
membership certificate for each use, exhibits its own finite subcollection.
For a PR checker in that formulation the supplied membership certificates
must themselves have a PR verification relation.

Each algorithm visits the nodes of the original $p$ once, carrying a
list of current maps or proofs downward. A current graft has the original
rule-tree size of its $q_i$ or is a one-node local assumption: weakening
and shifting relabel vertices. The grafted tree thus has at most
$|p|(1+\max_i |q_i|)$ rule nodes, taking the empty maximum to be zero.
During relabelling the number of nested extensions is at most $|p|$.
Contexts are finite lists, with at most $|p|$ new entries on a path.
Repeated applications of a fixed PR bound for a single relabelling
are bounded by its primitive iteration $|p|$ times. More formally one
can run a forward pass producing the local arguments at all addresses
and a reverse pass producing the result. The assertion that every
processed address has passed its local equality tests is a single
bounded Boolean fold. This supplies quantifier-free induction for the
certificate transformations.
\end{proof}

\begin{proposition}[A Hilbert deduction transformer]\label{prop:hilbert-ded}
Consider a named-variable Hilbert presentation with modus ponens,
generalization on a variable absent from the free variables of all
undischarged assumptions, and axiom instances
\[
 \begin{gathered}
  K: U\to(V\to U),\qquad
  S:(U\to(V\to W))\to((U\to V)\to(U\to W)),\\
  (\forall x(U\to V))\to(U\to\forall x V)
       \qquad(x\notin\operatorname{FV}(U)).
 \end{gathered}
\]
Other fixed logical axioms may be included. From an $N$-line derivation
of $B$ from $\Gamma,A$, a PR algorithm produces a derivation of $A\to B$
from $\Gamma$ in at most $5N$ lines.
\end{proposition}
\begin{proof}
For each old line $D$ store the location of its new last line $A\to D$.
For an axiom or an assumption in $\Gamma$, copy $D$, append
$D\to(A\to D)$ and apply modus ponens. For the hypothesis $A$ use the
two $K$ instances $A\to((A\to A)\to A)$ and $A\to(A\to A)$.
Apply the matching $S$ instance and modus ponens twice to obtain
$A\to A$ in five lines. For modus ponens on $E$ and $E\to D$ use $S$
and two applications of modus ponens to their stored versions. If $D$
is $\forall x E$ by generalization, $x$ is absent from $A$ and from
$\Gamma$. Generalize the stored $A\to E$, append the displayed
quantifier axiom and apply modus ponens. All new references point
backwards. Checking the stored-location invariant for the bounded
list of earlier lines proves correctness by induction on the scan.
\end{proof}

\subsection{Equality and eliminating explicit definitions}

\begin{proposition}[Derived equality calculus]\label{prop:derived-eq}
In $I$ there are uniform proofs of symmetry, transitivity, congruence for
each function and predicate, and replacement of equals in every formula.
Conversely reflexivity, symmetry, transitivity and congruence for
the function and predicate symbols suffice for formula replacement.
\end{proposition}
\begin{proof}
Given $s=t$, substitute into $z=s$, starting with $s=s$, to get $t=s$.
From $s=t$ and $t=u$ substitute in $s=z$ to get $s=u$. Insert a single
varying argument into the context
$f(\ldots,s,\ldots)=f(\ldots,z,\ldots)$ and start with reflexivity.
For a predicate the context is $R(\ldots,z,\ldots)$. Change
successive arguments one at a time and compose the implications or
identities. Replacement for an arbitrary context is precisely the
elimination rule, and symmetry gives the reverse implication.

For the converse fix $s=t$ and a term context $u(z)$. A variable occurrence
of the hole uses $s=t$. Another variable uses reflexivity. At a function
node apply its congruence to the already obtained identities of its
arguments. For formula contexts construct both directions together.
At a relation use term congruence and predicate congruence in each
direction. At $u=v$ use transitivity and the two term identities,
reversing one by symmetry. Both occurrences of $\bot$ are identical.
At $\wedge$ convert each component. At $\vee$ reason by cases. At an
implication, first convert its antecedent backwards, apply the given
implication and convert the result forwards. The reverse proof exchanges
these two roles. Under $\forall$, open at a fresh parameter, apply the
body conversion with the shifted $s=t$, and generalize again. Under
$\exists$ open a witness in a local subproof, convert its body and
introduce the existential with that very witness. The conclusion omits
the local parameter, as required by elimination. The procedure recurses
over the finitely many input nodes. Store both conversions per node
and compose them by the graft operation. An explicit PR size bound is
obtained by iterating a bound for this construction at one node the
number of nodes in the input. Here a step combines finitely many
previously constructed proofs, at most the maximum symbol arity plus
a fixed number of rule templates, and uses the PR operations of
Theorem~\ref{thm:graft}.
\end{proof}

\begin{theorem}[Conservative predicate abbreviations]\label{thm:def-abbrev}
Let $D$ be an old-language formula in $k$ slots. Extend the language
by one fresh $k$-ary predicate $P$ and the universally closed axiom
$\forall\vec x\,(P(\vec x)\leftrightarrow D(\vec x))$ with the indicated
argument order. Every extended proof of an old sequent from old axioms
and this axiom yields an old proof of the sequent by a PR transformation,
in $I$ or in $C$. A finite ordered list of definitions referring only to
earlier ones can likewise be eliminated.
\end{theorem}
\begin{proof}
Expand $P(\vec t)$ into $D[\vec t]$ at every occurrence in sequents
and formula annotations, and recurse through the old constructors.
Expansion commutes with substitution: at $P(\vec t)$ this is the
composition identity of Theorem~\ref{thm:subst}. Under a binder
the same identity uses the lifted substitution. Expansion is the
identity on the old language. It therefore preserves every rule
test, including opening in equality elimination. Replace the expanded
defining axiom by its old proof: derive $D\to D$ by discharge twice
for the two conjuncts of $D\leftrightarrow D$, combine them and
universally introduce each parameter. In an arbitrary ambient context
one may first prove this closed theorem, then weaken and shift as needed.
Use grafting to replace all uses of the axiom. For a list, process the
definitions in reverse order. Iteration over the input list of the
uniform single-definition algorithm is PR.
\end{proof}

\section{Soundness, quantifier laws and small separating models}
\label{sec:deduction-soundness}

Here is an interpretation that checks which logical principles the two
certificate types support. A finite Kripke frame consists of a nonempty
finite list of worlds, a reflexive transitive order $\leq$, and nonempty
finite domains $D_w$ of natural-number labels, with $D_w\subseteq D_v$ for
$w\leq v$. At each world there are total function tables and relation
tables on its domain. Require that functions preserve their values on
earlier tuples and positive relation entries persist upwards. Equality is
equality of the individual labels. All these conditions are decidable
by bounded scans. For an assignment $a\in D_w^n$ interpret terms by table
lookup. Denote its value by $[t]_{w,a}$. Atomic forcing uses the given
tables, $\bot$ has value false, and $\wedge,\vee$ are the Boolean
conjunction and disjunction at the current world. The remaining clauses
are
\[
\begin{aligned}
 w,a\Vdash U\to V
    &\iff (\forall v\ge w)\,(v,a\Vdash U\ \Rightarrow\ v,a\Vdash V),\\
 w,a\Vdash\forall B
    &\iff (\forall v\ge w)(\forall e\in D_v)\,v,(e,a)\Vdash B,\\
 w,a\Vdash\exists B
    &\iff (\exists e\in D_w)\,w,(e,a)\Vdash B .
\end{aligned}
\]
The displayed quantifiers range over finite tables. All three tests
return bits. With one world these are precisely ordinary two-valued
semantics, also called Tarskian semantics here.

\begin{lemma}[Persistence and semantic substitution]\label{lem:krip-sub}
Term values are unchanged at later worlds on old assignments, and positive
forcing persists. For a substitution $\sigma:n\to m$ and an assignment
$a$ in $D_w^m$, put $b_i=[\sigma_i]_{w,a}$. Then
\[
  [t[\sigma]]_{w,a}=[t]_{w,b},\qquad
  w,a\Vdash A[\sigma]\ \Longleftrightarrow\ w,b\Vdash A.
\]
Truth depends only on the free support. In particular a sentence has the
same truth value under every assignment at a fixed world.
\end{lemma}
\begin{proof}
Variable lookup is unchanged at a larger world. The condition on
function tables propagates this assertion through each term node. For
forcing, persistence at atomic nodes follows by the condition on relations
and stability of term values. The cases $\bot,\wedge,\vee$ follow from their
definitions. For an implication or a universal formula, every extension
of a later world is already an extension of the earlier one. For an
existential formula its witness still belongs to the larger domain
and its body persists by the child assertion.

For term substitution, variable lookup gives the claimed equality on
variables and applying the same function entry gives it on a compound
term. Atom evaluation then agrees. Conjunction, disjunction and falsity
preserve agreement. At implication, in each $v\ge w$ the values of
$\sigma$ at $a$ are still the original $b$, so apply agreement to the
two subformulas there. At a quantifier, the image of $\sigma^+$ evaluated
at $(e,a)$ is $(e,b)$: its first component is the fresh slot and the
rest ignore that slot. This works at each $v\ge w$ for the universal
and at $w$ for the existential. Finally assignments agreeing on free
support have equal term values at all corresponding occurrences.
At depth $d$ extend both by the same $d$ chosen individual labels
along a common sequence of worlds. Each variable lookup in the
corresponding bodies then agrees, and the forcing clauses give the
same output. This also proves the sentence assertion.
\end{proof}

\begin{theorem}[Finite soundness]\label{thm:krip-check}
Every accepted $I$ proof preserves forcing at every world of every
coded finite Kripke model. Every accepted $C$ proof preserves truth in
every coded nonempty finite two-valued structure. Both claims are
uniform theorems of iPRA.
\end{theorem}
\begin{proof}
For each node, tabulate for every world and every assignment the assertion
that if all its assumptions are forced then its conclusion is forced.
An assumption is a lookup. A false conclusion from an assumed proof of
$\bot$ poses no case, because a forced $\bot$ is the false bit. The
conjunction and disjunction rules read the corresponding bits. For
disjunction elimination choose the true disjunct and use its branch with
that assumption appended. For implication introduction pass to any $v\ge w$
forcing the antecedent. Persistence takes every old assumption to $v$,
so the branch gives the consequent. For elimination use the implication
test at $v=w$ and the true antecedent.

For universal introduction take arbitrary $v\ge w$ and $e\in D_v$.
Each old assumption, shifted, is true at $(v,(e,a))$ by persistence and
the semantic substitution lemma. The premise gives the body there.
For elimination choose $v=w$ and $e=[t]_{w,a}$, then use semantic
substitution. An existential introduction uses the same $e$ as witness.
For elimination scan for the first witness $e$ of the forced major premise
at $w$. At $(w,(e,a))$ the shifted old assumptions and the body all hold,
so the minor premise gives the shifted conclusion. Ignoring the fresh
slot gives the desired conclusion at $(w,a)$. For parameter erasure
use the least label of nonempty $D_w$ and the same shift identity.
Reflexivity uses identity of labels. In equality elimination, $s,t$
have the same value $e$, so both opened formulas have the truth of the
same body at $(e,a)$. In a one-world model implication is the usual
Boolean implication. Its double-negation operation is the identity
on each of the two bits, verifying the additional $C$ rule.

For an explicit induction invariant fix the input model and proof.
If the label universe has size $h$ after reindexing and the largest
arity at any local sequent or subformula is $k$, there are at most
$|W|\max(1,h)^k$ world-assignment rows per occurrence. Both $h$ and
$k$ come from finite scans, including extra slots under binders.
Compute truth and persistence in bottom-up syntax order, and soundness
in bottom-up proof order. At each stage take the bounded conjunction of
the assertions for all earlier rows. This is a quantifier-free induction
in the language of PR functions, even though the exposition speaks of
every assignment at once.
\end{proof}

\begin{corollary}[Consistency and explicit separation]\label{cor:ded-cons}
Neither calculus has a proof of $\bot$ from the empty context.
There is no $I$ proof of excluded middle $H\vee\neg H$ for a fresh
nullary predicate $H$, whereas $C$ has a proof.
\end{corollary}
\begin{proof}
Take the one-element one-world model with constant functions and empty
relation tables for consistency. For separation, take two worlds $r<s$
with a singleton domain and force $H$ only at $s$. At $r$ neither $H$
nor $\neg H$ holds. Finite soundness excludes an $I$ certificate.
In $I$, if $\neg(H\vee\neg H)$, then $H$ gives a contradiction by the
left injection, yielding $\neg H$. The right injection now gives a
contradiction too. Discharge to obtain $\neg\neg(H\vee\neg H)$ and
apply the extra rule of $C$.
\end{proof}

\subsection{Which quantifier identities have intuitionistic proofs?}

Write $P(x),Q(x)$ for formulas with an indicated slot and $H$ for a
formula independent of that slot. Other parameters may occur throughout.
The following equivalences have uniform $I$ proofs.
\begin{align*}
 \neg\exists x\,P(x)&\ \leftrightarrow\ \forall x\,\neg P(x),\tag{1}\\
 \forall x\forall y\,P(x,y)&\ \leftrightarrow\ \forall y\forall x\,P(x,y),\tag{2}\\
 \exists x\exists y\,P(x,y)&\ \leftrightarrow\ \exists y\exists x\,P(x,y),\tag{3}\\
 \forall x\,(P(x)\wedge Q(x))&\ \leftrightarrow\
      (\forall x\,P(x))\wedge(\forall x\,Q(x)),\tag{4}\\
 \exists x\,(P(x)\vee Q(x))&\ \leftrightarrow\
      (\exists x\,P(x))\vee(\exists x\,Q(x)),\tag{5}\\
 \exists x\,(H\wedge Q(x))&\ \leftrightarrow\ H\wedge\exists x\,Q(x),\tag{6}\\
 (H\to\forall x\,P(x))&\ \leftrightarrow\ \forall x\,(H\to P(x)),\tag{7}\\
 ((\exists x\,P(x))\to H)&\ \leftrightarrow\ \forall x\,(P(x)\to H).\tag{8}
\end{align*}
For (1), a putative $P(u)$ contradicts $\neg\exists x\,P(x)$ by
existential introduction. Generalize $u$. Conversely a witness for
$\exists x\,P(x)$ contradicts its particular negation. For (2), introduce
the target's two fresh parameters, instantiate the input in the opposite
order and generalize in the target order. For (3), open the two input
witnesses in order and introduce the same two in the reversed nesting.
Both operations reverse themselves on exchanging the roles of $x,y$.
For (4), specialize at a fresh parameter and project either component,
generalizing the two results separately. For the converse specialize
both inputs at the same parameter, pair them and generalize. For (5),
open the witness and split its disjunction, introducing the matching
existential and disjunct in each case. In reverse, split the outer
disjunction, open its witness and rebuild the inner disjunction
with that witness. For (6), retain $H$ and introduce $\exists x\,Q(x)$
from the witness, or, in reverse, carry $H$ into the witness branch
and pair it with $Q$. For (7), take a fresh $u$ and assume $H$.
Use the implication and specialize its universal result. In reverse,
assume $H$, specialize at fresh $u$, apply and generalize.
For (8), a given $P(u)$ supplies an existential witness to the left
implication. In reverse, open a given existential at $u$ and apply the
specialized right-hand universal there. Independence of $H$ makes the
parameter conditions in (6)--(8) hold. Each argument is a fixed proof
template, using shifts and openings on input codes.

There are also $I$ proofs of
$\exists x\,\neg P(x)\to\neg\forall x\,P(x)$ and
$H\vee\forall x\,P(x)\to\forall x\,(H\vee P(x))$.
For the first, specialize the supposed universal at the negative
witness. For the second split the disjunction. At each fresh individual
use the left injection if given $H$, and otherwise specialize the
universal and use the right injection.

\begin{proposition}[Two converses separated]\label{prop:quant-krip}
The converses of these last two implications have uniform $C$ proofs.
For a fresh unary $P$ and nullary $H$, neither converse has an $I$ proof.
\end{proposition}
\begin{proof}
In $C$, assume $\neg\forall x\,P(x)$ and, toward a contradiction,
$\neg\exists x\,\neg P(x)$. At any fresh $u$, $\neg P(u)$ gives a
forbidden witness, so $\neg\neg P(u)$. Eliminate double negation and
generalize, contradicting the first assumption. Discharge the second
assumption and eliminate double negation once more. For the other
converse assume $\forall x\,(H\vee P(x))$ and split on $H\vee\neg H$.
In the first branch inject $H$. In the second, arbitrary $u$ gives
$H\vee P(u)$. The left alternative contradicts $\neg H$, and the right
gives $P(u)$. Generalize and use the right injection.

For the first countermodel take a root $r$ with two incomparable
terminal successors $s,t$. All domains are $\{0,1\}$. Force no
$P$ at $r$, just $P(0)$ at $s$ and just $P(1)$ at $t$. The universal
fails at all three worlds, so its negation holds at $r$. Each potential
negative witness is defeated at one successor, so the existential
fails there at the root. For the second take $r<s$, with
$D_r=\{0\}$ and $D_s=\{0,1\}$. Force $P(0)$ throughout, never
$P(1)$, and force $H$ just at $s$. The premise holds at $r$: its
local choice is $P(0)$, and either individual at $s$ satisfies $H$.
At the root $H$ fails and $\forall x\,P(x)$ fails by the later
individual $1$. Apply Theorem~\ref{thm:krip-check} in both cases.
\end{proof}

\paragraph{Lean implementation path for the proof checker.}
Define raw inductives \texttt{Term} (variable, function on a list)
and \texttt{Form} with the constructors above, and Boolean functions
\texttt{termOK} and \texttt{formOK} taking the arity and signature.
One nondependent implementation of substitution carries the binder
depth explicitly and uses the depth algorithm of
Theorem~\ref{thm:subst}. Define \texttt{Sequent} with natural-number
arity, assumption list and conclusion, and a proof node with its sequent,
tag, annotation and ordered list of children. For termination one can
implement tree traversals over postorder arrays: each entry refers only
to earlier entries and the structural test verifies the child sorts and
arities. For a context-dependent calculation, first unroll occurrences
(a subtree repeated by two parents gets two addresses), propagate
arities or substitutions downward, then fold upwards. If the array length
is $N$, an unrolling has depth at most $N$ and branching at most the
maximum declared child-list length $b$. The expression $1+(b+1)+\cdots+(b+1)^N$
bounds its size. Ill-formed inputs return failure before traversal.
Define \texttt{localOK} by cases on the tags listed above, and
\texttt{check} by conjunction with the stored child bits. Prove
\texttt{lift\_comp} and \texttt{open\_subst} from the composition
law. Prove the node equalities for renaming and grafting by cases
on the rule tag, using these two lemmas at quantifiers and equality
elimination. Implement forcing using finite tables, with two extra
loops at a universal node and a bounded witness lookup at an existential
node. Prove the persistence and substitution tables before the
per-rule invariant in Theorem~\ref{thm:krip-check}. For an internal
iPRA implementation, replace inductives and arrays by their natural
codes, and generate the same bounded-fold identities. Lean induction on
the conjunction of all earlier Boolean checks then translates to
quantifier-free induction with the model, formula, proof and other
input codes as parameters.

The companion \texttt{NegativeTranslation.lean} implements the typed
rule trees and proves \texttt{FOLNegative.finite\_sound} and
\texttt{FOLNegative.consistency}. Its argument-list datatype permits
function symbols of arbitrary finite arity. Scope and the signature's
arity checks belong to the external syntax encoding. The Lean model
uses $\operatorname{Fin}(n+1)$, arbitrary total operations and decidable
relations. It also proves soundness for a presented nonempty carrier
when the semantic double-negation step is available for each enabled
classical rule. The finite expanding-domain Kripke argument above uses
the same quantifier rules with an additional world parameter.

\section{Interpreting classical proofs intuitionistically}\label{sec:negative}
The article distinguishes classical from intuitionistic first-order logic. There is nevertheless an arithmetical equivalence between their proof predicates after a change of formula. The following variant of the Gödel--Gentzen translation keeps the positive constructors under double negation.

\begin{theorem}[Negative translation on certificates]\label{thm:negative}
Let $A^N$ replace each atom (including equality) by its double negation, preserve $\bot,\wedge,\to,\forall$, and send $A\vee B$ and $\exists x A$ to $\neg\neg(A^N\vee B^N)$ and $\neg\neg\exists x A^N$. For every finite context,
\[
             \Gamma\vdash_C A\quad\Longleftrightarrow\quad
             \Gamma^N\vdash_I A^N .
\]
Both directions have PR maps on certificates. In particular the two empty theories are equiconsistent in iPRA.
\end{theorem}
\begin{proof}
Write $JX=\neg\neg X$. There are fixed intuitionistic derivations of $X\to JX$, of $JX\to JY$ given $X\to Y$ (compose a refutation of $Y$ with that implication), and of $JJX\to JX$ (given $k:\neg X$, a proof of $JX$ contradicts $k$). Say that $Y$ is stable when $JY\to Y$ is derived. For stable $Y$ an implication $X\to Y$ extends to $JX\to Y$ by the second derivation and stability.

Every $A^N$ is stable, with a uniform certificate. All double-negated heads use the third derivation. For $\bot$, apply the assumed $J\bot$ to the identity of $\bot$. For conjunction use the two projections and their stable conclusions. For implication $X\to Y$, under a hypothesis $x:X$ the map from $X\to Y$ to $Y$ is application to $x$. Extend it to $J(X\to Y)\to Y$ and discharge $x$. For a universal formula, fix a fresh individual $a$, use the map from $\forall x\,Y(x)$ to $Y(a)$, extend by stability at $a$, and generalize. Each premise comes from the corresponding child formula with any needed shift. Translation commutes with term substitution because it changes no term and introduces no new individual binder except the binder already present.

Translate the proof from its leaves up. Hypotheses, implication and conjunction rules, universal rules, explosion and parameter erasure translate directly. Disjunction introduction and existential introduction add $X\to JX$ after the corresponding ordinary rule. In a disjunction elimination the translated branches define an ordinary implication from the disjunction of translated disjuncts to the stable translated conclusion. Extend it to the double negation of that disjunction. For existential elimination use the translated fresh-variable branch to obtain an implication from $\exists x A^N$ to the same stable conclusion, then extend it. Equality introduction follows by double-negating reflexivity. For transport along $s=t$, assume the translated minor premise, perform ordinary transport in the formula context $B^N$ under the temporary assumption $s=t$, and discharge. Extend the resulting implication to $J(s=t)$ using stability of $B^N\langle t\rangle$. The translated classical elimination is exactly the stability derivation, since $(\neg\neg A)^N=J(A^N)$.

In classical deduction $A^N\leftrightarrow A$ has a PR proof: atom, disjunction and existential heads shed or introduce $J$ by classical elimination or intuitionistic introduction. Use the child equivalences under each constructor. For implication use the reverse equivalence on the antecedent by composition, for a quantifier instantiate a fresh variable and then generalize or eliminate its existential as appropriate. Start the reverse transformation with the given intuitionistic tree as a classical one, supply each translated hypothesis by its forward equivalence under $\Gamma$, graft these proofs, and use the reverse equivalence at the conclusion. Since $\bot^N=\bot$, the consistency claim follows by composing either transformation with a supposed contradiction.

All templates in this proof are fixed finite rule trees. A pass through the formula constructs its translation and its stability and equivalence certificates. A pass through the input proof constructs the translated proof from child certificates. Copies and substitutions have the bounds of Section~\ref{sec:deduction}. With $s$ the current maximum code, one stage is bounded by the finite maximum of the fixed template builders on inputs at most $s$ and the input labels. Iterate this bound once per row. Correctness is the conjunction of the local proof-checker tests on those rows.
\end{proof}

\section{A constructive completeness theorem with syntactic worlds}\label{sec:canonical}
There is a sense of completeness in which iPRA can reason about worlds of unbounded size. We give the exact semantics because a change in the interpretation of disjunction or of falsity would change the theorem. Start with the fragment consisting of atoms, equality, $\bot,\wedge,\to,\forall$. Worlds are pairs $w=(n,\Gamma)$ of an arity and a finite context in the full intuitionistic calculus $I$. An arrow $w\to v=(m,\Delta)$ consists of a term substitution $\sigma:n\to m$ and a list of proofs $\Delta\vdash_I G[\sigma]$, one per $G$ in $\Gamma$. Identity arrows are hypothesis leaves. For composition, substitute the second term map into each proof in the first list and graft the proof list of the second arrow. Theorem~\ref{thm:graft} verifies this operation and the transport of any derivation along an arrow. Multiple certificates can witness the same access between contexts.

For each \emph{fixed} formula $A$ with $k$ slots define an arithmetic predicate $w\Vdash A[\theta]$ with a numerical parameter $\theta:k\to n$. In this paragraph $\vdash$ denotes existence of an accepted proof code. For atoms (including equality) and $\bot$, use $\Gamma\vdash_I A[\theta]$. Conjunction uses conjunction in arithmetic. An implication $B\to C$ holds if, for every coded arrow $(\sigma,\pi):w\to v$, the assertion $v\Vdash B[\theta;\sigma]$ implies $v\Vdash C[\theta;\sigma]$. A universal $\forall B$ holds if for every such arrow and every $m$-ary term $t$ one has $v\Vdash B[(t,\theta;\sigma)]$. The tuple on the right sends slot zero to $t$ and old slots to the composite. These clauses are a finite compilation of the syntax of $A$ into an ordinary arithmetic formula, with potentially unbounded quantifiers over arrow codes. They require no decision of the resulting predicates.

\begin{theorem}[Schematic forcing equals derivability]\label{thm:syntacticforcing}
For each formula $A$ of this fragment, iPRA proves, uniformly in $w,\theta$,
\[
            w\Vdash A[\theta]\quad\longleftrightarrow\quad
            \Gamma\vdash_I A[\theta].
\]
Forcing is persistent under these arrows. For a sentence the equivalence at the empty context gives completeness for the indicated interpretation.
\end{theorem}
\begin{proof}
Construct a finite iPRA argument in the metatheory by induction on the fixed template $A$. This is an induction that \emph{generates proofs}, not an internal induction on formulas with a quantified induction invariant. The atomic and falsity cases use their definitions. For conjunction, project either part of a given proof, or combine the two proofs supplied by the induction hypotheses. These are fixed PR operations.

Given a proof of $B\to C$ at $w$ and an arbitrary arrow, transport it. The induction hypothesis converts forcing of $B$ at the target into a proof there. Apply the transported implication and use the other induction hypothesis to get forcing of $C$. In reverse, take $v=(n,\Gamma,B[\theta])$ with identity term map and the weakened context proofs. Its last hypothesis supplies forcing of $B$ by the induction hypothesis. Apply the assumed forcing implication at this arrow and turn its conclusion into a proof of $C[\theta]$ in $v$. Implication introduction returns the required proof in $w$.

A proof of $\forall B$ transports along each arrow and specializes to every term $t$. The substitution law gives exactly the body instance in the definition. Apply the body induction hypothesis. For the reverse take the single generic arrow to $(n+1,\uparrow\Gamma)$ whose term map is the shift, with the shifted hypotheses as its proof list. Specialize there at $v_0$. The combined body substitution is $(v_0,\uparrow\theta)$, the lift of $\theta$. The body induction hypothesis returns a derivation under $\uparrow\Gamma$. Universal introduction gives the claimed universal in $w$. Finally transport the proof on the right of the equivalence and use the equivalence again to obtain persistence. Every use of existential elimination on a proof predicate provides an actual certificate to the finite constructions involved.
\end{proof}

\begin{corollary}[An arithmetical completeness reading for classical theorems]
For each classical sentence $A$ one can compute a formula $A^-$ of the indicated fragment and an arithmetic sentence $\mathsf{Force}(A^-)$ such that iPRA proves
\[
       (\exists p\,\operatorname{Check}_C(p,0,[],A)=1)
                   \ \longleftrightarrow\ \mathsf{Force}(A^-).
\]
\end{corollary}
\begin{proof}
Use the translation of Theorem~\ref{thm:negative}, writing the disjunction clause as $\neg(\neg X\wedge\neg Y)$ and the existential clause as $\neg\forall x\neg X(x)$. Intuitionistically, a refutation of $X\vee Y$ is equivalent to a pair of refutations by the disjunction rules, and a refutation of $\exists x X(x)$ is equivalent to $\forall x\neg X(x)$ by the quantifier rules. Negating these equivalences gives the required replacements, including substitutions in enclosing formulas. Compose the proof maps and use the preceding theorem at $w=(0,[])$ with the empty substitution.
\end{proof}

At a world carrying a contradiction, falsity is forced and every formula of the fragment is forced. Equality here is provable equality and can grow along an arrow. These facts distinguish the result from Tarski completeness and from the usual nonexploding Kripke semantics with rigid identity. The formula compiler produces an iPRA theorem separately for every $A$. No single arithmetical satisfaction predicate for arbitrary arithmetic formulas is asserted.

\section{Finite structures and the truth definition}\label{sec:models}
A signature is a finite list of function and relation symbols with their arities. A constant is a function of arity zero. Write $[m]=\{0,\ldots,m-1\}$, with $m>0$. A structure $M$ of size $m$ gives each $s$-ary function a table of $m^s$ entries in $[m]$ and each $s$-ary relation a bit table of the same length. The lexicographic index of a tuple is obtained by iterating $j\mapsto mj+a$ over its entries. The empty tuple has index zero. Fixed lengths make this work even at $m=1$. Code validity, including all range checks, is a bounded scan. Statements involving mathematical structures below are uniform implications about valid codes.

\subsection{Evaluation and change of environment}
An environment for $k$ free slots is a tuple $e\in[m]^k$. Write $a::e$ for extension at the front. For de Bruijn indices, $v_0$ under a new binder names its new entry. Write $t^M(e)$ for term evaluation and $[A]_M(e)$ for the truth bit. A variable reads an entry and a function node indexes its table by the values of its sons. A relation node does the same with a bit table, while $t=u$ tests numerical equality. Falsity gives $0$, conjunction takes minimum, disjunction takes maximum, and implication takes $\max(1-p,q)$ on input bits $p,q$. Quantifiers give the finite folds
\[
 [\forall A]_M(e)=\min_{a<m}[A]_M(a::e),\qquad
 [\exists A]_M(e)=\max_{a<m}[A]_M(a::e).
\]
One can use $1$ and $0$, respectively, as the initial values of these folds. Negation is implication to falsity. These are classical \emph{object-language} truth tables evaluated by constructive arithmetic.

\begin{proposition}[Uniform model checking and witnesses]\label{prop:mc}
Term values and formula truth are primitive recursive (PR) uniformly in the signature, structure, expression and environment. A true existential supplies its least witness and a false universal supplies its least counterexample. A closed sentence has a well-defined truth bit independent of any extra environment entries.
\end{proposition}
\begin{proof}
List the nodes of the syntax tree in postorder. For each occurrence at binder depth $d$, allocate one cell for each of the $m^{k+d}$ environments of that length. Fill rows in node order, applying the displayed operations to cells in earlier rows. If the tree has $L$ nodes, at most $L m^{k+L}$ cells are needed, and each quantifier cell scans $m$ entries. A node's variable, symbol and argument indices and the table arithmetic all use bounded operations on the input or on previously constructed rows. Encoding each updated row into a single number gives ordinary numerical primitive recursion. Return zero on malformed input. Induction on the filled rows proves their defining equations. For the witness, scan the child row until the first $1$ (respectively $0$), returning sentinel $m$ on failure. Induction on a scan prefix says that the sentinel persists exactly while all entries read have the opposite value. Environment independence follows from the next lemma.
\end{proof}

\begin{lemma}[Coincidence, renaming and semantic substitution]\label{lem:semsub}
Environments agreeing on the free slots of an expression give it the same value. For a capture-avoiding simultaneous substitution $\sigma=(s_0,\ldots,s_{k-1})$, let $e_\sigma=(s_0^M(e),\ldots,s_{k-1}^M(e))$. Then
\[
 (t[\sigma])^M(e)=t^M(e_\sigma),\qquad
 [A[\sigma]]_M(e)=[A]_M(e_\sigma).
\]
Consequently, equal-valued terms can replace each other in any formula context.
\end{lemma}
\begin{proof}
For coincidence, variable lookup uses a slot on which the environments agree. Equal lists of child values select the same function or relation entry and give the same equality or Boolean bit. Under a quantifier extend both environments by the same $a$. They agree at the new slot and on the shifted free slots needed by the body. The folds therefore agree entry by entry. The same traversal proves the renaming formula: replacing $v_i$ by $v_{\rho(i)}$ evaluates an expression at the tuple $(e_{\rho(i)})_i$. Under a binder the index map fixes $0$ and sends $i+1$ to $\rho(i)+1$, so the two ways to form each extended tuple coincide.

For substitution, the variable case reads precisely $s_i^M(e)$. All function and atomic cases follow by equality of their child values. At a binder the lifted tuple sends its first slot to $v_0$ and old slot $i+1$ to $s_i$ with all its free indices raised by one. By renaming, this raised term evaluates on $a::e$ to $s_i^M(e)$. The environment built from the lifted tuple is thus $a::e_\sigma$. Apply the body identity to these environments and take the same minimum or maximum. Boolean nodes use equal input bits. For replacement, take the formula with a free placeholder at the indicated occurrences and apply the identity twice. The two resulting environments agree. All of these inductions can take the conjunction of the required equalities at the nodes already visited, over all assignments of the row lengths just specified. For substitution the lifted tuple at each node is predetermined by its depth. It is a computed parameter of that node. Thus no induction hypothesis quantifies over an unbounded set of substitutions or environments.
\end{proof}

\begin{corollary}[Closure, duality and finite evaluation games]
The universal closure of $A$ is true in $M$ exactly when all assignments satisfy it. Its existential closure is true exactly when some assignment satisfies it. Negating a quantifier exchanges it with its dual and negates the body. In the ordinary evaluation game for a formula in negation normal form, the verifying player wins exactly when its truth bit is $1$. A winning strategy for the appropriate player is PR in the input.
\end{corollary}
\begin{proof}
Repeated quantifier folds enumerate the Cartesian power, in the order set by the binders. Complement interchanges min and max: scan the same entries on the two sides, starting with complementary bits. At a disjunction or existential the verifier chooses the branch or element. At a conjunction or universal the falsifier chooses. A literal awards the game to the verifier exactly at bit $1$. At a verifier choice with bit $1$, choose the first child of bit $1$. At a falsifier choice with bit $1$ every child has bit $1$. The symmetric rule keeps bit $0$ for the falsifier. Each move goes down the syntax tree, so the invariant persists to an awarded leaf. Storing the choices with the evaluation cells produces both conditional strategies by bounded scans. Negations can be pushed into atoms by the bit identities just proved and the two-element Boolean truth tables.
\end{proof}

\subsection{Maps, equality and finite diagrams}
A map $h:[m]\to[n]$ is a length-$m$ table. A homomorphism $M\to N$ commutes with every function table and carries every true relation entry to a true entry. An embedding is injective, commutes with functions and preserves and reflects relations. An isomorphism is a surjective embedding. Each condition is a finite test. A substructure has a nonempty subset of the domain closed under all functions (including constants), with each relation restricted to that subset. Enumerating its members in increasing order puts it in our coding convention.

\begin{theorem}[Preservation under structure maps]\label{thm:maps}
Homomorphisms preserve existential-positive formulas, built from atoms, truth, falsity, conjunction, disjunction and existence. Embeddings preserve and reflect quantifier-free formulas. Existential sentences with quantifier-free matrix pass from a substructure to the larger structure. Universal ones pass to a substructure. Isomorphisms preserve and reflect all formulas with corresponding parameters.
\end{theorem}
\begin{proof}
Compute first $h(t^M(e))=t^N(h e)$, where $h e$ means applying $h$ entrywise. The variable case is lookup and the function step is the commutation equation. For a homomorphism this proves preservation of equality and reduces preservation of atomic relations to its definition. Conjunction carries both successes, disjunction carries its successful branch, and existence transports its least witness $a$ to $h(a)$. Truth and falsity have constant bits. For an embedding injectivity reflects equality and the defining relation equivalence gives equality of the remaining atomic bits. Every Boolean connective preserves equality of bits. Witnesses for an existential block in a substructure therefore work in the large structure. A universal block true in the large structure gives a true matrix on each tuple from the substructure, which reflects back. For an isomorphism scan for the unique preimage of each element. At quantifiers every extended environment in one structure has a corresponding extension in the other. The body bits coincide, so do their extrema. These are inductions on the finite evaluation tables. In particular, a parameter-free definable relation is invariant under every automorphism, and the same holds with parameters under each automorphism fixing them.
\end{proof}

\begin{theorem}[Equality from a congruence]\label{thm:eqquot}
Let $E$ be a binary table on $M$ satisfying reflexivity, symmetry and transitivity. Suppose functions respect $E$ coordinatewise and relations have equal bits on coordinatewise $E$-equivalent tuples. There is a PR quotient $Q$ and surjection $p$ such that, for every formula $A$ of the original signature with logical equality,
\[
 [A^{E}]_M(e)=[A]_Q(p e).
\]
Here $A^{E}$ replaces each logical equality atom $t=u$ by $E(t,u)$.
\end{theorem}
\begin{proof}
Set $r(a)$ to the least $b<m$ with $E(a,b)$. Reflexivity ensures success. If $E(a,b)$ then symmetry and transitivity give equal sets of candidates, hence $r(a)=r(b)$. Conversely, if $r(a)=r(b)$ then $E(a,r(a))$ and $E(b,r(b))$ give $E(a,b)$. Thus $r(r(a))=r(a)$. Enumerate the fixed points, and let $p(a)$ be the ordinal position of $r(a)$ in this list. For each quotient tuple apply the original function on the representatives and then apply $p$. Read predicates directly on representatives. The assumed congruences give the same answers on all equivalent tuples, so induction on terms identifies quotient values with projections of their old values. Equality of projected values is exactly $E$ of old values. The relation case uses saturation. Boolean steps preserve equal bits. Every quotient value has its listed representative, and each original value has its projection, so the extrema defining either quantifier agree. Every test scans at most a finite table or a finite power of the carrier. In particular this realizes the equality-axiom reduction on codes for finite structures.
\end{proof}

\begin{theorem}[Two finite diagrams]\label{thm:diagrams}
Given a finite $M$, one can compute an existential-positive sentence $P_M$ and a sentence $D_M$ of quantifier rank at most $m+1$ such that for every finite $N$ in the same signature
\[
 [P_M]_N=1\ \Longleftrightarrow\ M\to N\text{ by a homomorphism},\qquad
 [D_M]_N=1\ \Longleftrightarrow\ M\cong N .
\]
In particular, agreement on all sentences between two finite structures implies isomorphism. Existence of a homomorphism is equivalent to inclusion of their existential-positive sentence theories in the same direction.
\end{theorem}
\begin{proof}
Name each $i<m$ by a variable $x_i$. For each function row include $f(x_{i_1},\ldots,x_{i_s})=x_{f^M(\bar i)}$. For each true relation row include $R(x_{i_1},\ldots,x_{i_s})$. Conjoin these clauses and existentially bind all names, using truth for the empty conjunction. This is $P_M$. A witnessing tuple is exactly the table of a homomorphism. If all existential-positive truths transfer, $P_M$ transfers because the identity witnesses it in $M$. The converse is Theorem~\ref{thm:maps}.

For $D_M$, insert under the same existential prefix the negative literal for every false relation row, $x_i\ne x_j$ for all $i<j$, and the clause $\forall y\,\bigvee_{i<m} y=x_i$. A witnessing tuple now enumerates the target without repeats and preserves and reflects every atomic table. A given isomorphism supplies that tuple. There are $m$ outer quantifiers and at most one further nested quantifier. This also proves finite categoricity from elementary equivalence via the single sentence $D_M$. In both constructions finite truth provides witnesses by repeated least search. The function and relation clauses include nullary rows.
\end{proof}

\subsection{Three changes of language}
\begin{proposition}[Graphs of functions]\label{prop:fgraph}
Functions and constants can be eliminated in favor of relation symbols in finite satisfiability, with a uniform PR translation and model recovery.
\end{proposition}
\begin{proof}
Introduce a fresh $(s+1)$-ary $G_f$ for each $s$-ary $f$. Require
$\forall\bar x\,\exists y\,G_f(\bar x,y)$ and
$\forall\bar x\,\forall y\,\forall z\,((G_f(\bar x,y)\land G_f(\bar x,z))\to y=z)$.
For a term $t$ and fresh $z$ construct its value formula $V_t(z)$. Put $V_{v_i}(z)=(z=v_i)$. For $t=f(t_1,\ldots,t_s)$ put
\[
 V_t(z)=\exists y_1\cdots y_s
 \left(\bigwedge_{j=1}^s V_{t_j}(y_j)\land G_f(y_1,\ldots,y_s,z)\right).
\]
In the nullary case use $G_f(z)$. Replace each atomic formula on terms by existential binding of their values, conjunction of the value formulas and the same atom on the fresh variables. Translate connectives and quantifiers compositionally, shifting free indices whenever extending scope. A structure with functions supplies their graphs. In the reverse direction search the output column of each input row. Totality supplies a first output and uniqueness identifies every output with it. Term induction gives exactly one $z$ satisfying $V_t(z)$, the original term value, hence every translated atom has the old bit. Boolean and quantifier folds now agree. Conjoining the graph axioms with the translated sentence proves the claim. Each constructor expands a finite list of previously computed children by a bounded fold.
\end{proof}

\begin{proposition}[Many sorts by relativization]
For a finite nonempty list of inhabited sorts, finite many-sorted satisfaction translates uniformly to one-sorted satisfaction. In the relational encoding, a finite model of the translated sort and typing axioms yields a finite many-sorted model.
\end{proposition}
\begin{proof}
Use the graph translation within each sort first if functions occur, choosing witnesses of the stated output sort. Make a tagged disjoint union of the carriers and add predicates $S_i$ for the tags. Assert $\exists x\,S_i(x)$ for each $i$, that every element belongs to some $S_i$, and that no element belongs to two. For each relation require that every true tuple has the declared coordinate sorts. Quantification over sort $i$ becomes $\exists x(S_i(x)\land A)$ or $\forall x(S_i(x)\to A)$. Copy equalities and predicates on well-typed variables. The filtered min/max runs over precisely the elements of the intended sort, since out-of-sort entries contribute the neutral bit. This proves correctness on every well-typed environment by the evaluation recurrence. Conversely enumerate each tag predicate in increasing order and restrict relation tables to the typed Cartesian products. The partition identifies the disjoint union of these enumerations with the given one-sorted carrier. The relativized graph axioms recover function tables if needed.
\end{proof}

\begin{proposition}[Relational algebra for finite assignments]
In a relational language the sets of satisfying $k$-tuples of first-order formulas are exactly the relations computed from the named tables, equality and full domain powers by coordinate substitutions, Boolean operations and one-coordinate projections. The translations in both directions are PR.
\end{proposition}
\begin{proof}
Interpret a formula in $k$ slots as its bit table on $[m]^k$. An atom $R(v_{i_1},\ldots,v_{i_s})$ pulls back a named table along the indicated coordinate list, with repeated coordinates permitted. Equality pulls back the diagonal table. Truth gives the full power, negation takes complement within that power and conjunction takes intersection. Existential quantification drops the first coordinate by maximum over its fibers. Other connectives and universal quantification follow by truth tables and duality. Conversely a named table is an atom on its coordinates, the diagonal is equality and the full power is truth. Pullback is capture-free variable renaming, Boolean operations are connectives and projection binds the coordinate being removed, after permuting it to slot zero. Each recursion reconstructs the same bit by its defining equation. The availability of full domain powers specifies the complement convention even if an element never occurs in a true relation row.
\end{proof}

\section{Finite back-and-forth and a sharp rank bound}\label{sec:ef}
Work in a finite relational signature with equality and without function symbols or constants. Nullary relations are allowed. With constants one may first use the translation in Proposition~\ref{prop:fgraph}, at a potential cost in quantifier rank. The quantifier rank $q(A)$ is zero on atoms, is the maximum of the child ranks at a Boolean connective, and increases by one at a quantifier.

\subsection{A game that returns a formula}
For $M,N$ of sizes $m,n$, a position consists of equal-length tuples $u,v$ of previously paired elements, ordered most recent first. Define $I(u,v)$ to say that all equalities $u_i=u_j$ and all $R^M(u_{i_1},\ldots,u_{i_s})$ have the same bits on $v$. Thus $I$ checks a finite list, including the empty index tuple for each nullary relation. Write $W_0=I$ and define $W_{r+1}(u,v)$ as the conjunction of $I(u,v)$ and
\[
 \begin{split}
 &\min_{a<m}\max_{b<n} W_r(a::u,b::v)=1,\\[-2pt]
 &\min_{b<n}\max_{a<m} W_r(a::u,b::v)=1 .
 \end{split}
\]
Treat predicates throughout as their bits. In one round Spoiler chooses one side and one element, and Duplicator chooses an element from the other side. Duplicator must maintain the partial isomorphism $I$. The first successful inner entry above provides a strategy reply. If a clause fails, its first failing outer entry is a Spoiler move for which every reply loses one level lower. Recursing on the number of rounds proves finite determinacy with strategies for either outcome. Testing $I$ at intermediate rounds is equivalent to testing it only at the end: a failing atomic comparison persists after any tuple extension, and the carriers are inhabited.

\begin{theorem}[Finite Ehrenfeucht--Fra\"iss\'e theorem with a separator]\label{thm:ef}
If $W_r(u,v)=1$, every formula in the indicated slots of quantifier rank at most $r$ has equal bits at $(M,u)$ and $(N,v)$. If $W_r(u,v)=0$, a PR procedure returns a formula of rank at most $r$ with value $1$ at $(M,u)$ and $0$ at $(N,v)$. Thus agreement on all formulas of this rank is equivalent to winning the game.
\end{theorem}
\begin{proof}
For agreement, follow the formula with a remaining budget $s$ initially equal to $r$. At an atom use $I$. Boolean operations preserve equal child bits, at the same budget. At an existential with a witness on the left, the first direction of $W_s$ pairs it with a right witness at budget $s-1$ by the body assertion. A right witness is transported back by the second direction. A universal is handled by transporting counterexamples in the same way. Budget suffices because each quantifier descent removes one from the rank. This proof uses induction on syntax, including Boolean nodes when the budget is zero.

For the other direction orient every separating formula to be true on the left. If $I$ fails, take the first discrepant atom or its negation. Otherwise the budget is positive. A failed first direction supplies $a$ such that each $b<n$ has a separator $B_b$ for $(a::u,b::v)$ at budget $s-1$. Then
\[
                  \exists x\,\bigwedge_{b<n} B_b(x,\bar z)
\]
separates: the one left witness $a$ makes all conjuncts true. For each proposed right witness $b$ its own conjunct fails. If the second direction fails, take $b$ such that each $a<m$ has a separator $C_a$ for $(a::u,b::v)$. Then
\[
                  \forall x\,\bigvee_{a<m} C_a(x,\bar z)
\]
is true on the left because each $a$ has its corresponding true disjunct, and false on the right at the chosen $b$ because every disjunct is false. Child formulas are already expressed in the new slot and shifted old slots. Neither finite connective increases rank, giving the rank bound.

For a numerical implementation, put $k=|u|$. The row at budget $i$ covers tuples of length $k+r-i$. There are at most $(r+1)(mn)^{k+r}$ cells across these rows. Compute the rows in increasing $i$ and store the bit, strategy choice and, on failure, the oriented formula. If $s_0$ is the maximum of two and the relation arities, a separating syntax tree need have at most $H_i$ nodes, with for example
\[
 H_0=s_0+5,\qquad H_{i+1}=(m+n+1)(H_i+4)+5 .
\]
Here indices count as single variable nodes. Their numerical labels are bounded by $k+r$, and relation labels by the signature. Binary folds realize the displayed finite connectives. Encoding the successive finite rows is ordinary PR even when formula codes exceed their node counts. The loop invariant checks the computed separator's rank and its two evaluation bits for every failed cell filled, and for each successful cell verifies its finite response clauses. Its induction formula is a bounded Boolean test. For the agreement half fix a supplied formula and conjoin the required implications over its nodes, all tuple pairs of the needed lengths and budgets at most $r$, proceeding in postorder. The equivalence with agreement on all formulas follows by testing the computed separator if the game fails.
\end{proof}

\begin{corollary}[A finite rank test for isomorphism]
For finite relational structures, $W_{\max(m,n)}((),())=1$ exactly when $M\cong N$.
\end{corollary}
\begin{proof}
An isomorphism and its inverse specify a reply to every move. Conversely play all elements on the larger side in enumeration order. A winning strategy pairs distinct elements with distinct elements, so the two finite sizes must agree. To see the cardinal inequality arithmetically, mark each new image in a bit vector. Its count increases by one at every move and never exceeds the length of the vector. At equal sizes every vector entry is marked. The final pairing is therefore a bijection, and $I$ checks every relation row on this enumeration. All moves and marks belong to bounded loops.
\end{proof}

\subsection{What rank can count on finite linear orders}
Let $L_m$ be $[m]$ with its usual strict order as the only nonlogical symbol. Write $T_0=0$, $T_{r+1}=2T_r+1$, so $T_r=2^r-1$.

\begin{theorem}[Exact threshold and nondefinability of parity]\label{thm:orders}
$L_m$ and $L_n$ agree on all sentences of rank at most $r$ exactly when $m=n$ or $\min(m,n)\ge T_r$. No first-order sentence in the single-order signature defines even cardinality on all nonempty finite linear orders. Given any proposed sentence one computes a concrete adjacent even--odd pair on which it gives the same bit.
\end{theorem}
\begin{proof}
Sort the distinct matched points. Maintain their equality and order, and require that the numbers of unplayed points in corresponding gaps (including the end gaps) are equal or both at least $T_s$, where $s$ rounds remain. Repeating a point is answered by its match. Given a new point in an equal-sized gap, copy its offset from the left. In a pair of unequal gaps each has size at least $2T_{s-1}+1$. If the chosen point leaves fewer than $T_{s-1}$ elements on its left, copy that left count. The two right counts are then at least $T_{s-1}$. If it leaves fewer on its right, copy the right count. If neither is small, choose a point with exactly $T_{s-1}$ elements on its left in the responding gap: both its sides are large enough. The argument works for moves on either side and also when $T_{s-1}=0$. The invariant at zero rounds ensures all equalities and order bits agree. At the initial position its condition is exactly the asserted sufficient size condition.

For sharpness, for each $0\le h\le T_s$ construct a formula saying that an interval has at least $h$ elements, of rank at most $s$, with its existing endpoints as free parameters. A missing endpoint imposes no bound. At $h=0$ take truth. At $h>0$ put $p=\min(h-1,T_{s-1})$, $q=h-1-p$, so $p,q\le T_{s-1}$. Ask for a point $x$ inside the interval, with at least $p$ elements in the interval to its left and $q$ in the one to its right. The interval guard uses only atomic order tests. Formally each child uses the appropriate original endpoint and $x$ as its other endpoint, under a fresh binder. If the interval has at least $h$ elements, its $(p+1)$st point is a witness. Conversely, the two disjoint open subintervals and the chosen point provide at least $p+q+1$ elements. Induction verifies the claim by bounded checks of all interval endpoints and $h\le T_s$ at each stage (allow endpoint sentinels and empty intervals). In orders of unequal sizes with $\min(m,n)<T_r$, choose $h=\min(m,n)+1$ and use the whole interval. The resulting sentence separates. Finally, given a sentence of rank $r$, take $m=2^{r+1}$ and $n=m+1$. The first is even, the second odd, and the sufficient condition forces agreement. This gives the uniform arithmetical counterexample claimed.
\end{proof}

\paragraph{Formalization details for the model and game algorithms.}
Use natural-number codes with operations for constant-length tuple enumeration, table lookup, prepending, and indexed left folds. Give \texttt{evalTerm} and \texttt{truth} their postorder implementations. The data lemmas needed for the binder proofs are \texttt{lookupConsZero}, \texttt{lookupConsSucc}, the identity between lookup on a mapped tuple and mapping a lookup, and the minimum/maximum witness identities for a finite bit row. Prove \texttt{evalRename} first, then \texttt{evalLift} and \texttt{evalSubst}. Formulate each as a rowwise equality guarded by code validity. The structure-map proofs reuse this pattern with a fixed map table. For quotients, implement \texttt{rep} by bounded minimum, prove $E(a,b)\leftrightarrow r(a)=r(b)$, then prove term projection before formula projection. For diagrams and language changes, assemble the displayed finite lists of clauses by indexed folds. Apply the closure and substitution lemmas to their evaluated matrices. For the game, store \texttt{win}, \texttt{choice} and \texttt{separator} per position. Prove the bounded row invariant in Theorem~\ref{thm:ef} and separately the implication from a winning position to agreement on each input formula. For the order corollary implement sorted pebbles and gaps by scans and the two-sided reply by the three cases in the proof. The companion \texttt{ef\_certificates.py} constructs and evaluates actual separating formulas, including pointed examples, and tests the rank threshold on small orders.

\section{Witness closure and the Tarski--Vaught test}\label{sec:witnessclosure}
The downward Löwenheim--Skolem construction closes a seed under functions selecting witnesses. Here is its finite version and a version for a presented arithmetic domain. The latter specifies exactly what an effective witness selection must supply.

\begin{theorem}[Finite elementary substructure for a specified fragment]\label{thm:TV}
Let $M$ be a nonempty finite structure of size $m$, $S$ a nonempty list of seed elements, and $\Phi$ a finite subformula-closed list in the Boolean basis $\bot,\neg,\wedge,\exists$, with atomic equality and relations. There is a PR substructure $N\subseteq M$ containing $S$ such that for each $A\in\Phi$ and each appropriately long tuple from $N$ one has $[A]_N=[A]_M$. The construction returns all the witness selections. For the least-witness policy below it returns the least subset containing $S$ and closed under those selections and the basic functions.
\end{theorem}
\begin{proof}
For each $\exists x B(x,\bar y)$ in the list and tuple $\bar a$ in $M$, compute $g_B(\bar a)$ as the least witness if one exists and the first seed otherwise. Thus all auxiliary functions are total, including those for sentences with an empty parameter tuple. Starting with $S$, add in parallel all values of these functions and all original function values on tuples from the current set, including all nullary functions. Use a bit mask. The process is increasing. Every nonstationary stage adds at least one new element. The count of marked elements at stage $i$ is at least $\min(m,|S_{\rm distinct}|+i)$ until stationarity. Therefore $m$ iterations reach a fixed point. Every closed superset contains each stage by induction, giving the minimality claim. On the final set restrict relations and functions, and list its members in increasing order for the code of $N$.

Term values computed in $N$ are the same elements as in $M$ by function closure. Equality and relation bits consequently agree. Boolean operations preserve agreement. At $\exists x B$, a witness in $N$ transfers to $M$ by the body result. If the formula holds in $M$ at a tuple from $N$, its selected witness is in $N$ by closure, and the body transfers back. Tabulating the bounded assignment rows gives the quantifier-free induction. Conversely, for any given substructure, agreement on the fragment forces the ordinary witness condition: when $M$ satisfies an existential at a tuple of the substructure, the substructure has some witness making the body true there. With inductive agreement on the body this is precisely the Tarski--Vaught test restricted to $\Phi$.
\end{proof}

\begin{proposition}[A presented Skolem hull]\label{prop:hull}
Fix a PR domain predicate $D(a)$, a seed $a_0$ in it, PR operations closed on $D$, and PR atomic relation bits. Fix a finite fragment as above. Supply for its subformulas PR truth bits with the semantic clauses on $D$ proved in iPRA, and for every true existential bit supply a PR function selecting a witness in $D$ with true body bit. A PR term presentation gives a substructure of this domain, generated by the seed, that agrees on the fragment. The evaluation map of this presentation and its equality test are PR.
\end{proposition}
\begin{proof}
Adjoin auxiliary function symbols for these selectors, replacing their values off a true existential by $a_0$, and a name for $a_0$. Use closed terms as names. Evaluate them by postorder recursion. Every value is in $D$ by closure. Two names are equivalent if their numerical values coincide. For an actual domain of unique codes, replace each valid term code $t$ by the least well-formed ground term code $u\le t$ of equal value. This is a bounded search, and two terms of equal value have the same representative: if $t\le t'$ then $t$ is already a candidate for $t'$ and all candidates smaller than the minimum for $t$ fail in both searches. The minimum is idempotent. The representatives constitute a decidable domain. Apply an original function by forming its term on representative arguments and taking the new representative. Read relations on the numerical values of the names. Evaluation now gives an injective homomorphism reflecting relations. Its range contains $a_0$ and is closed under each selector by the same constructor operation. Membership in that range inside $D$ need not be decided.

For each fixed existential in $\Phi$, a true ambient bit on a tuple of names gives a witness name by applying its selector symbol. A witness in the presentation has an ambient value, and the body induction transfers in either direction. Atomic and Boolean clauses follow as before. This gives, for each fixed subformula, equality between its ambient semantic reading and its interpretation using quantifiers over representative codes. Equivalently the ambient bit pulled back along term evaluation supplies its PR truth bit on the hull. The hypotheses already specify the truth and witnesses on the larger domain, so the uniform new operation is building the smaller presentation. Cardinal arithmetic and a choice operation on arbitrary structures play no role in this statement.
\end{proof}

\section{Normal forms, witnesses and refutations}
\label{sec:normal}

This section connects three kinds of evidence. A finite truth table can supply witnesses for a Skolem expansion. A ground truth table can supply a resolution refutation. A cut-free derivation can supply the very ground instances to which the second construction applies. All transformations below take numerical codes as input. Statements involving evaluation in a structure use a coded nonempty finite domain $[d]=\{0,\ldots,d-1\}$ unless expressly stated otherwise.

\subsection{Prenex and clausal normal forms}

Use $\neg,\wedge,\vee,\to,\forall,\exists$, together with atomic equality, relation atoms and truth constants. Treat $A\leftrightarrow B$ as an abbreviation for $(A\to B)\wedge(B\to A)$. Formula sizes below count the expanded trees, including term nodes. In a traversal, give each binder a distinct temporary name greater than every free identifier. For de Bruijn input, the name of a bound occurrence is found by looking up its binder in the traversal stack. A reverse traversal with a scope stack returns de Bruijn syntax. Thus a prefix can be moved without ever changing which binder owns an occurrence.

\begin{theorem}[Polarity and prenex conversion]
There are primitive-recursive transformations $N^+,N^-,P$. The first two yield formulas in negation normal form and evaluate respectively to the original bit and its complement. $P$ yields a string of quantifiers followed by a quantifier-free matrix with the same free variables and the same evaluation as the input.
\end{theorem}
\begin{proof}
A positive atom stays an atom and a negative atom becomes its negation. A negation exchanges the two polarities. On a conjunction the positive branch joins the positive children by $\wedge$, and the negative branch joins the negative children by $\vee$. On a disjunction interchange $\wedge$ and $\vee$ in those instructions. For implication the positive output is $N^-(A)\vee N^+(B)$ and the negative output is $N^+(A)\wedge N^-(B)$. The positive quantifier branch keeps the quantifier, and the negative branch interchanges $\forall,\exists$ and takes the negative body. Complement truth constants at negative polarity. The two outputs have the stated bits at an atom. The four possibilities for a pair of child bits prove each binary step. Complementing a finite conjunction of bits gives the disjunction of their complements: the empty folds have bits $1,0$, and appending one bit applies the binary De Morgan law. The dual identity is identical with the two bits exchanged. These identities prove the quantifier steps.

Run a second traversal on the positive output. A literal returns the empty prefix and itself. A quantifier prepends its binder to the child's prefix. Given prefix--matrix pairs $(U,B),(V,C)$ for two children joined by $\star\in\{\wedge,\vee\}$, return $(UV,B\star C)$. For the proof, move the binders of $U$ across the right child and then those of $V$ across the left matrix, from outside to inside. At each move one operand has a fixed bit $b$ independent of the variable moved. Both finite folds commute with $z\mapsto z\star b$: if this map is the identity this is immediate, and if it is constant, a fold of $d\ge1$ copies of that constant is the same constant. Distinct fresh names give exactly the independence required. This also explains the domain hypothesis: on the empty domain $\top\vee\exists x\,\bot$ and $\exists x(\top\vee\bot)$ have different bits.

Postorder processing constructs at most two polarity outputs per node. A single polarity output visits each input occurrence once and adds at most one negation per atomic occurrence. Prenex conversion moves binder records and retains the matrix nodes. The temporary-name range has length the number of binders. For the induction verifying conversion, fix the input and structure, and conjoin the two desired bit identities over the processed subtrees and all assignments to names appearing there. If $v$ names appear, there are at most $d^v$ assignments. Each induction assertion is one primitive-recursive bit test.
\end{proof}

The same fold calculation gives the article's identities $\neg\forall x A\equiv\exists x\neg A$ and $\neg\exists x A\equiv\forall x\neg A$ in these models. Two like quantifiers commute: a double conjunction is $1$ exactly when every ordered-pair entry is $1$. For two disjunctions, scan for the first $1$ and transpose its pair of indices. Universal quantification distributes over conjunction and existential quantification over disjunction by distributing the corresponding finite folds entry by entry. Freshness is essential in the rule for pulling a quantifier across an operand. The preceding proof records the point where it is used.

\begin{theorem}[Skolem tables with an explicit choice function]
Let $F=Q_1x_1\cdots Q_qx_q B$ be a closed prenex sentence with quantifier-free $B$. For every existential position $i$ add a fresh function $s_i$ of arity the number $u_i$ of preceding universal variables, in their prefix order. Substitute $s_i(\bar x_{<i}^{\forall})$ for $x_i$ in $B$, keep the universal prefix, and call the result $S$. A primitive-recursive function expands every finite model of $F$ to a model of $S$ on the same domain. Every finite model of $S$ reduces to a model of $F$.
\end{theorem}
\begin{proof}
For each prefix history $h$ of length $i$, compute the bit $T_i(h)$ of the remaining suffix. Fill the last row with the matrix bits, then compute the preceding rows by finite AND or OR according to the next quantifier. Now fill the function tables in prefix order. At an existential position and a fixed tuple of preceding universal moves, reconstruct the intervening existential moves using tables already filled. This gives one history $h$. Record the least $a<d$ for which the next row reads $T_i(h,a)=1$ (here $i$ is the one-based position just played). Use $0$ when there is no such entry. This fills the entire table even when the root bit is zero. If the root bit is one, a universal move preserves a one for each choice and an existential table selects a one whenever its parent is one. Induction on the play length therefore makes the final matrix one for every universal tuple. Evaluation after substitution gives exactly this matrix value.

Conversely, suppose all universal tuples give one under some supplied Skolem tables. Every history consistent with the tables has all of its complete continuations true when later existential moves also follow the tables. Work backward from length $q$. At a universal position apply the preceding assertion to every child. At an existential position the table entry gives one child to which it applies. The suffix bit is one in both cases, hence the root bit is one. Each row assertion ranges over at most $d^q$ histories. The suffix array has $\sum_{j=0}^q d^j$ entries, and the new function tables use $\sum_{i:Q_i=\exists}d^{u_i}$ entries. All bounds include the constant-function case $d^0=1$. Each replacement of a variable occurrence adds at most $q$ argument nodes. These counts give primitive recursions on prefix positions and bounded scans for every operation.
\end{proof}

\begin{proposition}[Two clause conversions]
A quantifier-free matrix in negation normal form has an equivalent conjunction of clauses. If it has $m$ nodes, one may use at most $2^m$ clauses of width at most $m$. Alternatively, after adding fresh relation symbols, at most $3m+1$ clauses suffice and every satisfying interpretation of the matrix at all universal tuples has a unique expansion satisfying all these clauses.
\end{proposition}
\begin{proof}
Use one singleton clause for a literal, no clauses for $\top$, and the single empty clause for $\bot$. Concatenate clause lists at a conjunction. At a disjunction form every pairwise union of one clause from each child list. If both children have a false clause, their union is false. A bounded scan finds such clauses whenever both lists have Boolean value zero. If either list is true, every union is true. This proves the disjunction rule including empty lists. Conjunction adds list lengths, disjunction multiplies them, and a union's width is bounded by the sum of widths. Induction gives the two bounds.

For the second construction let $\bar u$ be all the universal variables. Use the literal itself as its output wire. For each binary node introduce a distinct predicate $p(\bar u)$. If its children's wires are $a,b$, impose the following clauses, where a bar complements a literal:
\[
\begin{array}{c|lll}
p\leftrightarrow(a\wedge b)&\bar p\vee a&\bar p\vee b&p\vee\bar a\vee\bar b\\
p\leftrightarrow(a\vee b)&p\vee\bar a&p\vee\bar b&\bar p\vee a\vee b .
\end{array}
\]
A constant node uses a fresh wire fixed by one unit clause, and a final unit asserts the root wire. For each tuple the three clauses force exactly the displayed Boolean gate value, as checking the two possible output bits shows. Consequently postorder evaluation supplies the only possible new predicate entries. The root clause holds for all tuples exactly when the matrix does. Arity zero covers a sentence with no universal variables. The bound counts clause occurrences, not the total length of their argument lists.
\end{proof}

\subsection{Ground resolution and analytic Herbrand witnesses}

A clause is a finite set of signed atoms. Repetition is immaterial. A list of clauses is true when every clause contains a true literal. Its empty clause is always false. A ground atom is an atom whose terms have no variables. In propositional tests distinct ground atoms are treated as independent bits until explicit equality or congruence conditions are added.

\begin{theorem}[Bounded resolution completeness and soundness]
For any finite family $G$ of ground clauses on $k$ distinct atoms, primitive-recursive functions return either a satisfying assignment or a resolution refutation with at most $2^{k+1}-1$ nodes. Each leaf names an input clause. Each internal node names a pivot and the union of its parents after deleting their designated opposite pivot literals.
\end{theorem}
\begin{proof}
Enumerate the $2^k$ assignments and return the least satisfying one if it exists. Otherwise at every full assignment choose the first input clause it falsifies. For a partial assignment $s$ to the first $j$ atoms let $B_s$ contain the $j$ signed literals falsified by those choices. At a leaf the selected clause is a subset of $B_s$. Work up the complete binary tree, retaining at each node a derived clause $D\subseteq B_s$. If $p$ is its next atom, the children have clauses $D_0\subseteq B_s\cup\{p\}$ and $D_1\subseteq B_s\cup\{\neg p\}$. If either lacks its displayed pivot, retain that child alone. Otherwise resolve the children on $p$. The union of the two remainders is a subset of $B_s$. At the root this forces the empty clause. Pruning children only decreases the given node bound. The induction assertion is the conjunction of the subset tests in one tree level.

For soundness fix a valuation satisfying both parents of a resolution step. If the pivot bit is zero, the parent containing its positive occurrence must have a true remainder. If the bit is one, the parent containing its negative occurrence must have a true remainder. In either case the union is true. Induction over a proof therefore precludes an empty conclusion under any satisfying valuation. Both the checker and the witness construction are bounded traversals.
\end{proof}

For universally closed clauses, an additional sound rule instantiates a clause by any displayed well-scoped substitution. Its truth follows by interpreting the substituted terms, then using the assignment so obtained in the original clause. A first-order resolution certificate may supply substitutions on two separately renamed clauses and a check that the resulting pivot atoms coincide. It may also factor two literals which that substitution makes equal. Literal-set union implements the latter deletion. These operations require only syntactic comparison. Soundness follows from the preceding rule and the ground pivot argument applied at each assignment. In particular a finite, propositionally inconsistent family of ground instances refutes its universal parent clauses. For example, $R(a)$, $\neg R(x)\vee R(f(x))$ and $\neg R(f(f(a)))$ have a four-leaf certificate. Use instances $x=a,f(a)$ of the middle clause and resolve successively to obtain $R(f(a))$, $R(f(f(a)))$, and the empty clause.

\begin{theorem}[Two-way cut-free Herbrand extraction]
\label{thm:herbrand}
In equality-free classical one-sided sequent calculus, take a sequent of closed formulas $E_i=\exists\bar x_i\,A_i$, with quantifier-free matrices in negation normal form. From an explicit cut-free tree proof one can extract a finite list of ground instances of the $A_i$ whose disjunction is a propositional tautology. Conversely, such a list and a successful tautology check give a cut-free proof of the original sequent. Both maps are primitive recursive. If the language has no closed term, use a fresh parameter $c$ in the ground language.
\end{theorem}
\begin{proof}
Use multisets (or lists with exchange), weakening, contraction, the one-sided rules for conjunction, disjunction, and existential instantiation, and initial sequents containing complementary literals or $\top$. The conjunction rule has two premises with the same context. The disjunction rule has one premise containing both children. The existential rule takes $\Gamma,A[t/x]$ to $\Gamma,\exists x A$. These are precisely the logical rules needed on paths read backward from this endsequent: weakening removes a formula, contraction duplicates one, and the logical rules peel its prefix or its matrix. Equivalently, restrict an analytic calculus in negation normal form to this endsequent. Atomic initial sequents and explicit structural rules fix the proof format.

Replace all free proof parameters by a fixed closed term (adjoin $c$ if necessary), after avoiding capture. Trace each root occurrence backward with a tag recording its root index and witnesses already supplied in prefix order. Copy tags at contraction and at branching. Whenever a prefix is exhausted, record its ground matrix and tuple, including occurrences whose prefix was empty at the root. Further decomposition of its matrix requires no new record. There are at most the number of root occurrences plus the number of existential rule nodes many records.

Fix any bit assignment on the recorded atoms, extended by zero on any other proof atoms. Suppose every recorded matrix is false. Give every still-prefixed occurrence the temporary value false. Give all quantifier-free occurrences their usual propositional value. Starting from the root, whose occurrences are all false, follow a path with every occurrence false. At an existential step the exposed occurrence still has a nonempty prefix or is a just-recorded false matrix. At a disjunction both children are false. At a conjunction choose a false child and follow its premise. Backward weakening, contraction and exchange preserve the assertion. No initial sequent can have all occurrences false. This contradicts the finite path ending at a leaf. Since the purported failure was a bit assignment, a loop over all assignments proves the tautology test succeeds. This argument follows only bounded paths. It also works when the recorded list is empty, in which case no proof could have supplied it.

Conversely start with the quantifier-free sequent consisting of the listed instances. Decompose a disjunction into its two disjuncts in one premise, or a conjunction into its two alternative premises. Propositional validity of the conclusion implies validity of each resulting premise by the two gate truth tables. Continue until $\top$ closes a branch or only literals and $\bot$ remain. A valid leaf without $\top$ must contain opposite literals, since otherwise assigning each occurring positive atom zero and each occurring negative atom one falsifies the whole leaf. This constructs a cut-free derivation using at most $2^{L+1}-1$ tree positions for $L$ connective occurrences on the initial list: along a branch the count strictly decreases. Now restore the recorded existential binders around each matrix in reverse prefix order using the recorded terms. Contract repeated copies of each $E_i$ and weaken in any missing $E_i$. If $c$ was adjoined and was absent from the intended signature, replace it throughout the proof by a single fresh free proof variable. No free variable enters the endsequent.

For numerical verification, keep tags and instance lists on a bounded traversal of the supplied tree. The invariant for extraction is indexed by proof nodes and valuations of the finite set of proof atoms. Grounding the free parameters and instantiating displayed terms are syntax traversals. The reverse construction recurses on the connective count with the displayed exponential bound. This gives quantifier-free induction for the advertised maps.
\end{proof}

Taking $A_i$ to be the conjunction of complements of a universal clause gives a direct corollary. Cut-free refutability of a finite equality-free universal clause family is equivalent to the existence of an explicit finite list of ground instances having a resolution refutation. The forward map uses Theorem~\ref{thm:herbrand} and bounded resolution. The reverse map uses resolution soundness, finite truth-table duality and the reverse construction in Theorem~\ref{thm:herbrand}. A witness list supplies its own term-depth bound. No bound on search depth as a function of the clause family is invoked.

\subsection{A decision with two verifiable outputs}

\begin{theorem}[Bernays--Sch\"onfinkel--Ramsey with equality]
\label{thm:bsr}
Let $F=\exists x_1\cdots x_e\,\forall y_1\cdots y_u\,B$ be closed. The finite signature has relation symbols, equality, $b$ constants, and no other functions. There is a primitive-recursive decision algorithm which returns either a model of $F$ of size at most $h=\max(1,b+e)$, or a checkable classical first-order derivation of $F\to\bot$. In particular the same bit decides whether $F$ has \emph{any} coded finite model and detects infeasibility even when equality can identify different existential witnesses or constants.
\end{theorem}
\begin{proof}
Use a name for each original constant and each existential witness. Use one dummy name only if there are no names. There are $h$ names in total. Let $G$ be the propositional CNF obtained by the following instructions.
\begin{enumerate}
\item Give the CNF conversion of $B$ the witness names for $\bar x$ and instantiate $\bar y$ in all $h^u$ ways by names. Include every resulting clause.
\item Represent each equality of names $i,j$ by an atom $E_{ij}$, and each relation on name tuples by an atom $R_{\bar i}$. Include the units $E_{ii}$, all clauses expressing $E_{ij}\to E_{ji}$, and all clauses expressing $E_{ij}\wedge E_{jk}\to E_{ik}$.
\item For each relation $R$ of arity $r$ and each pair of name tuples include the clause for
\[
  \left(\bigwedge_{t=1}^r E_{i_t j_t}\right)\wedge R_{\bar i}
  \ \longrightarrow\ R_{\bar j}.
\]
\end{enumerate}
There are at most $k=h^2+\sum_R h^{\operatorname{arity}(R)}$ atoms. In addition to the ground matrix clauses, the construction uses at most $h+h^2+h^3+\sum_R h^{2\operatorname{arity}(R)}$ equality and relation clauses. Apply bounded resolution.

From a satisfying assignment, form the equivalence classes of $E$ using as representative the least equivalent name. Reflexivity guarantees a successful bounded search. Symmetry and transitivity imply two names have the same representative exactly when $E$ relates them. Interpret names by their representatives, and assign $R$ on representative tuples using the assigned bits. The congruence clauses in both tuple orders imply that replacing every coordinate by its representative preserves its bit. Thus for every tuple of name representatives the matrix evaluates to one. Every representative is itself a name, so all universal assignments are covered. The designated witness names supply the original existential moves. Ordering the representatives and replacing each by its rank gives a genuine model on an initial segment of size between one and $h$.

For the opposite branch, a resolution leaf has one of two justifications. For an equality or congruence leaf use logical equality: repeatedly replace one tuple coordinate using its equality, thereby obtaining the displayed congruence implication. Reflexivity is an axiom, symmetry follows by replacing in $i=i$, and transitivity by replacing in $i=j$. For a matrix leaf, assume $F$, open its existential witnesses using fresh parameters, and apply universal elimination at the indicated tuple of names. The inferred ground matrix entails each of its CNF clauses. To obtain classical proof codes uniformly for this last step, split on each of the finitely many atoms by excluded middle. Under a complete assignment, formulas with value one have proofs, and formulas with value zero have proofs of their negation, from the signed atomic assumptions: recurse over the formula, using the connective rules and the known signs of its children. In the true-matrix case the indicated clause has value one. In the false-matrix case the assumed matrix gives a contradiction. Disjunction elimination on the successive excluded middles discharges the signed assumptions. The same method gives the small propositional conversions from the equality implications to clauses. A resolution step is derivable by case analysis on both parent disjunctions: if either supplies a remainder literal introduce it into the conclusion. The opposite pivot pair gives a contradiction and hence the conclusion. Repetition and order in a clause can be changed by the same disjunction rules. The empty-clause root therefore gives $\bot$. Discharge the fresh witness parameters by existential elimination in reverse order, then discharge $F$. The dummy name may be translated to a fresh free variable, as in Theorem~\ref{thm:herbrand}. Classical rules used here occur in the \emph{output proof}. Its generation and checking use bounded operations in iPRA.

Finally, any given finite model of $F$ supplies satisfying bits for $G$. Read its constants, scan for an existential tuple making the universal suffix true, and interpret the dummy name by element zero when needed. Record the equalities and relations of these elements. Every matrix instance is true, and actual equality supplies all congruence clauses. Hence the resolution branch is impossible for this model. The three constructions show exhaustiveness and exclusivity constructively. In ordinary external semantics the same finite restriction applies to any inhabited model with actual equality by taking the interpretations of these finitely many names. The internal result already decides between a finite model and a formal refutation.
\end{proof}

\paragraph{An equality test that changes the answer.}
For $F=\exists x\exists y\,(x=y\wedge P(x)\wedge\neg P(y))$ the ground matrix forces $E_{12},P_1,\neg P_2$. The congruence clause $\neg E_{12}\vee\neg P_1\vee P_2$ resolves against these three units to give the empty clause. For $\exists x\exists y\,(x=y\wedge P(x))$ both witnesses may name the same element and the quotient returns a one-element model. Assigning distinct domain elements to fresh constants by default would lose that possibility.

\paragraph{Lean implementation order.}
Use \texttt{Nat} identifiers and carry the prefix as a list of pairs (quantifier tag, fresh variable identifier). Define \texttt{polar : Bool $\to$ Formula $\to$ Formula}, \texttt{splitPrefix} returning the pair (prefix, matrix), and \texttt{suffixRows} by primitive recursion on the reversed prefix. Prove the bit equations for finite folds before proving correctness of \texttt{polar} and \texttt{splitPrefix}. A Skolem-table row takes a preceding universal tuple to the least accepted child of the reconstructed history. Its forward invariant is truth at every history generated from that tuple. Its reverse invariant is truth of the original suffix whenever all consistent completions satisfy the matrix. For resolution use signed integers or pairs (atom index, sign), a sorted deduplicated list for a clause, and an inductive certificate with constructors for indexed input leaves and binary pivots. The induction lemma returns both an accepted certificate and a clause contained in the literals falsified by the current prefix. For Herbrand extraction store a root index and a list of witnesses on each sequent occurrence, copying that record at contraction. Formalize the false-premise selector as a function of a valued sequent, its last rule and the valuation. Its iteration is bounded by tree height. For the fragment decision, build the lists of ground relation and equality atoms by tuple enumeration and prove representative equality in both directions before the structural induction on the ground matrix. Each correctness lemma is transferred to code arithmetic by collecting checks over the finite stage, assignments, and tuples specified in its proof.

\section{Monadic types, sharp bounds and graph reachability}\label{sec:expressive}
\subsection{Counting colours in a unary signature}
Consider $p$ unary predicates, $z$ nullary predicates, equality, and no constants or positive-arity functions. A colour $c<2^p$ is the $p$-bit truth vector of a point. Let $m_M(c)$ count its occurrences in a nonempty finite structure $M$. For $q\ge0$ put $h^M_q(c)=\min(q,m_M(c))$.

\begin{theorem}[A complete monadic invariant]\label{thm:monadic}
Two such finite structures agree on every sentence of quantifier rank at most $q$ exactly when their nullary bits and all their $h_q(c)$ agree. From a failed comparison one can produce a distinguishing sentence of rank at most $q$. Finite satisfiability and finite validity in this language are PR decidable. Every finitely satisfiable sentence of rank $q\ge1$ has a model of at most $q2^p$ points, a bound attained for every $q\ge1,p\ge0$. At rank zero a one-point model suffices.
\end{theorem}
\begin{proof}
In the $q$-round game, match previously chosen points to their partners. For a new point choose an unused point of the same colour. With $r$ rounds remaining maintain equality of the tuples' equality patterns and colours, and, for each colour, equality of unused counts truncated at $r$. If the chosen colour has a new point on one side, the other has one because $r>0$. When the counts below $r$ coincide they both decrease by one, and when both are at least $r$ the new counts are both at least $r-1$. Repeating a point only lowers the threshold. At the leaves all relation atoms and equalities on the tuple agree. This is the game of Section~\ref{sec:ef}. Bounded search selects the least eligible reply. Alternatively the same argument transfers an existential witness and a universal counterexample at each quantifier.

For the converse, a mismatching nullary atom distinguishes immediately. If the smaller truncated count is $k$ and the larger one exceeds $k$, let $D_c(x)$ be the conjunction prescribing all $p$ bits, using $x=x$ for an empty conjunction. Then
\[
 \exists x_0\cdots\exists x_k\,
  \Big(\bigwedge_{i\le k}D_c(x_i)\ \wedge\
                  \bigwedge_{i<j\le k}x_i\ne x_j\Big)
\]
has rank $k+1\le q$ and holds exactly on the larger side. This follows by sorting or scanning for distinct elements of that colour. If an orientation is required, negate it. Hence equality of all bounded-rank truths, although expressed with an unbounded formula-code quantifier, has a PR complete test with a PR witness to failure.

For $q>0$ retain the first $h_q(c)$ elements of every colour, then reindex. The retained structure is nonempty and has the same invariant. To decide, enumerate the $2^z(q+1)^{2^p}$ possible bit/count profiles and discard the all-zero count vector. Realize each by its listed number of points of each colour and evaluate the sentence. For $q=0$ take each nullary valuation and one point of arbitrary colour. To see sharpness, conjoin for all $c$ the assertion that at least $q$ distinct points have colour $c$, using a separate block of $q$ existential binders inside each conjunct. Its rank is $q$. Unequal colours cannot share a point. A model with $q$ points of each colour realizes the sentence at the bound. In the equality-free language, matching a colour permits unlimited reuse of the same partner. Already one point per realized colour suffices at every positive rank.
\end{proof}

\paragraph{A normal form for monadic sentences.}
For fixed $q>0$ the conjunction describing a profile specifies the nullaries and, for each colour of count $k<q$, asserts at least $k$ and fewer than $k+1$ points. For count $q$ it asserts at least $q$. The zero threshold means truth. All parts have rank at most $q$. Disjoin these conjunctions over the accepting profiles. By the theorem this Boolean combination of threshold assertions is equivalent to the original sentence on every finite input. The empty disjunction means falsity. At rank zero use the truth table of nullaries alone. Classically the same decision answers unrestricted monadic satisfiability, because an arbitrary structure admits the same truncation of finitely many colour classes. For an infinite presented domain, deciding whether an individual colour occurs need not be effective. The internal uniform algorithm acts on sentences and finite tables.

\begin{theorem}[Proof-producing monadic decision]\label{thm:monadiccert}
For a closed monadic sentence $A$ in the preceding signature, a PR procedure returns either a finite model of $A$ or a classical deduction of $\neg A$. Applying it to $\neg A$ decides classical theoremhood in this fragment with a proof or a finite countermodel.
\end{theorem}
\begin{proof}
For each colour $c$, write $H_{c,j}$ for the closed threshold formula just used (with $j$ witnesses, and $H_{c,0}=\top$). We describe quantifier elimination to a Boolean combination of nullaries, tests on free variables, and the $H_{c,j}$ treated as closed units. In addition to the term for the result the transformation constructs a classical proof of its equivalence to the input. Atom and Boolean steps use the corresponding logical rules. Replace a universal by $\neg\exists\neg$, with its classical proof.

It remains to eliminate $\exists x\,D$ after the child has been translated. Expand $D$ into disjunctive normal form treating a threshold as one propositional atom. The propositional equivalence has a proof by splitting on its finitely many atoms. At a full valuation conjunction, disjunction, implication and negation are introduced or refuted from their signed children. Distribute existence over the disjunction and work on one conjunction. First orient equality and inequality involving $x$ by symmetry. Remove occurrences of $x=x$ and discard any conjunctive case containing $x\ne x$. If some conjunct requires $x=y_i$, substitute $y_i$ throughout, remove the quantifier, and retain the resulting conjunction. In the forward direction eliminate the witness and transport along the required equality. In the reverse direction use $y_i$ as witness. If no such equality occurs, all the remaining conditions on $x$ prescribe or prohibit unary bits, and require inequality from a list of parameters.

Split on every colour of the remaining parameters, and every equality between them, using excluded middle on the corresponding atoms. Inconsistent patterns give contradictions by equality transport. Split also on the possible colour $c$ of the witness. These alternatives are exhaustive by the $p$ atomic cases. Discard colours inconsistent with the unary conditions on $x$. In a surviving case let $d$ be the number of distinct forbidden parameters of colour $c$, calculated from the equality pattern. Replace the existence of $x$ by $H_{c,d+1}$, retaining the case conditions and all literals independent of $x$. Indeed a witness and one representative of each forbidden equality class supply $d+1$ pairwise distinct points of colour $c$. Conversely, open $d+1$ distinct witnesses for the threshold. Compare them to the $d$ forbidden representatives. If all were forbidden, their matching indices would be $d+1$ distinct numbers below $d$, impossible by the elementary marking argument. The comparisons give a surviving witness by a finite decision tree of equality cases. Its colour already ensures inequality from the forbidden parameters of other colours. This gives an actual classical proof template at each finite $d$, generated by expanding those cases, including $d=0$.

For a closed input, when eliminating a binder under $h$ other binders its remaining free parameters number at most $h$, so every new threshold is at most $h+1\le q$. Closed threshold units introduced earlier retain their scope. At the root, decide all nullaries and all thresholds up to $q$ by classical case splits inside the desired proof. Derive $H_{c,j+1}\to H_{c,j}$ by omitting the last witness, which removes all nonmonotone bit patterns. When $q>0$ and all $H_{c,1}$ are denied, take a fresh individual, determine its colour and introduce it as a witness to contradict the corresponding denial. Parameter erasure closes that branch. All remaining profiles have the canonical models of Theorem~\ref{thm:monadic}. At $q=0$ use one point. If some model accepts, output it. Otherwise the root Boolean formula is false in each consistent profile, and its signed atomic assumptions prove its negation by propositional evaluation. Recombine the cases and use the compiled equivalence to conclude $\neg A$. Running on $\neg A$ and eliminating double negation gives the theoremhood assertion by finite soundness.

All DNF, pattern and witness-comparison stages enumerate finite bitstrings with lengths computed from the current expression, and each stage constructs a finite rule tree. Treat introduced thresholds as marked units while processing parents, expanding them only in the output proofs. A fixed PR majorant for one such stage can be iterated once per original syntax node. Structural correctness collects the local proof-check bits as in Lemma~\ref{lem:toolkit}.
\end{proof}

\subsection{Two finite graphs a sentence cannot distinguish}
Connectedness and reachability are explicit examples in the starting article. Both already fail to be definable over finite undirected graphs. A direct construction here avoids any appeal to compactness.

\begin{theorem}[Rank-bounded witnesses against reachability]\label{thm:cycles}
Given $q$, put $s=2^q$ and $L=8(q+3)(s+1)$. Let $G$ be a cycle of length $2L$, with distinguished vertices $0,L$. Let $H$ be two disjoint cycles of length $L$, with one distinguished vertex in each. These pointed graphs agree at the displayed pair on every formula of rank at most $q$ whose free variables lie in those two slots, in the language of adjacency and equality. The distinguished vertices are joined by a path in $G$ and lie in different components of $H$. Even without distinguished vertices $G$ and $H$ agree through rank $q$.
\end{theorem}
\begin{proof}
Orient each cycle in its cyclic coordinate system, solely to describe Duplicator's replies. At $r$ rounds remaining preserve, for every pair of marked vertices, every displacement $d$ with $|d|\le 2^r$: the second is the first plus $d$ in its cycle on one side exactly when this holds on the other. Cross-component displacements are always false. These intervals give unique displacements whenever one exists, by the lower bound on the cycle lengths. The two initial marks satisfy the invariant including displacement zero.

Put $t=2^{r-1}$ for a move with $r>0$. If the new vertex is an offset $d$, $|d|\le t$, from an old mark, return the same offset from the corresponding mark. Any offset $e$ of length at most $t$ between the new vertex and another mark gives an offset of length at most $2t$ between the two old marks. Its truth transfers in both directions and subtraction of $d$ gives exactly the desired offset. If the vertex is far from all old marks, choose on the other side the least vertex outside all their radius-$t$ intervals. There are at most $q+2$ old marks, each excluding at most $2t+1$ vertices. Their total is bounded by $(q+2)(2t+1)<2L$. A bounded scan therefore succeeds. All old incidences retain the narrower invariant, and the new vertex has just its zero self-offset in the range. The argument works whichever board was chosen. At $r=0$ the displacements $0,1,-1$ give equality and undirected adjacency. Rank invariance now gives the assertion. The same strategy ignoring an unused pair of initial marks gives the sentence version. Traversing $L$ positive steps in $G$ gives the path. In $H$ an induction on path length preserves the component label. A proposed defining formula codes its own rank and hence its own counterexample pair.
\end{proof}

\section{Computations encoded by first-order sentences}\label{sec:computability}
Two reductions answer different questions. The first converts a successful run into a deduction, and an unsuccessful \emph{halting} run into a finite countermodel. The second converts an arbitrary accepting run into a finite model and recovers a run from every such model. Both reductions have arithmetical converses or separation consequences that are provable without deciding whether an arbitrary run halts.

\subsection{The numerical machine model}
Use deterministic machines with two stacks over a finite alphabet. An instruction either pushes a fixed letter on a selected stack and jumps, or inspects and pops a selected stack, with one jump for empty (leave it empty) and a designated jump for each popped letter. There are two stopping labels, yes and no. The input word starts on the right stack, its first letter uppermost, and the left stack starts empty. A program is a numbered finite table. A missing instruction can be treated as a permanent loop. The bounded-run predicate iterates the table lookup and the list updates for $t$ steps, freezing a stopped configuration. It is PR. Program codes and input words convert to each other by any fixed self-delimiting binary serialization.

Here are the coding facts needed for the diagonal arguments. A two-stack machine simulates a one-tape machine by keeping the cells left of the head, nearest first, on the left stack and the current cell and cells to its right on the right. An empty read returns blank. To move right, pop the old current cell and push its new value on the left. To move left, pop the old current cell, push its new value on the right, then transfer the nearest left cell (or a blank) onto the right. A stay is a right-stack top replacement. A finite instruction block performs each update and can remember any one letter and transition in its control. Conversely, ordinary finite scans, copying and the evaluation of any \emph{fixed} PR definition can be programmed on a tape. Store a list of binary registers separated by delimiters. A scan copies or replaces a register using a marked cursor, and successor is ordinary carry. To implement composition reserve disjoint scratch registers for the component computations. To implement primitive recursion copy the counter, compute the initial value and run the step block once for each counter decrement, passing the iteration index, previous value and untouched parameters. Induction on this fixed program construction gives PR running-time bounds for each PR routine: the recursion case iterates a PR upper bound for running time and register sizes together. Its iPRA invariant records the contents of the designated registers at block boundaries. A tape with auxiliary tracks uses a fixed product alphabet. These constructions justify precomposing a hypothetical machine with each particular compiler below, making a program that calls it on a transformed input, swaps its two stopping labels, or runs an unbounded enumeration one stage at a time. Only its finite traces are used in iPRA.

\subsection{A proof or a finite countermodel}
\begin{theorem}[Constructive Church separation]\label{thm:church}
There is a PR map $(P,w)\mapsto F_{P,w}$ to closed sentences. From a yes trace obtain an intuitionistic proof of $F_{P,w}$. From a no trace obtain a finite model of $\neg F_{P,w}$. No total computable bit test accepts every intuitionistically provable sentence and rejects every sentence with a finite Boolean countermodel. This holds already with equality and just one nonlogical binary relation.
\end{theorem}
\begin{proof}
First use a constant $e$, a unary constructor $c_a$ for each letter, a binary predicate $C_i$ for each program label, and a nullary predicate $Y$. The word term $[u]$ is $e$ for empty and $c_a([v])$ for $u=av$. Form a finite list $T$ of Horn sentences, including $C_{\rm start}(e,[w])$ and $\forall l,r\,(C_{\rm yes}(l,r)\to Y)$. For push-left from $i$ to $j$ add
\[
             C_i(l,r)\to C_j(c_a(l),r).
\]
For pop-left add $C_i(e,r)\to C_{j_0}(e,r)$ and
$C_i(c_a(l),r)\to C_{j_a}(l,r)$ for each letter $a$. Close all clauses universally. Right-stack rules exchange the two argument positions. Stop labels have no instruction clauses. Put $F=(\bigwedge T)\to Y$. Given an accepting trace, start with the initial fact, instantiate the clause for each instruction with the concrete word terms, and use modus ponens. The final clause yields $Y$. Conjunction projections and a single discharge produce the required proof.

For a rejecting trace let $B$ be the maximum length of any stack that occurs in it. Take all words of length at most $B$ and an overflow point $*$. Interpret $e$ as empty, $c_a$ as prefixing $a$ if the result fits, and as $*$ otherwise, including on $*$. Set $C_i$ to exactly the pairs of actual stack words visited at $i$. Make $Y$ false. Every push from a true antecedent is a step of the trace and fits the bound. The only preimage of a true nonempty word $av$ under $c_a$ is $v$. It has no preimage under another constructor. Thus every matching pop clause from a true antecedent produces the actual next pair. No antecedent at yes is true. These finite tables satisfy $T$ and falsify $F$, also when $B=0$.

Suppose a separator $D$ is given. Make a program $J$ that on input the code $u$ computes $F_{u,u}$ and runs $D$ on it, going to no on bit one and yes on bit zero. Take $u$ to be the code of $J$. Totality supplies the stopping trace. On bit zero the first construction supplies a proof that forces $D$ to output one. On bit one the second supplies a finite countermodel that forces zero. This contradiction is intuitionistic: the disjunction is read off a supplied terminating trace. Equivalently the arithmetic theorem says that no machine with a proof of total two-valued behavior can satisfy the two universally stated separator conditions, and, internally, negates the existence of a code satisfying their conjunction. The argument uses the test's claimed totality as a logical hypothesis.
\end{proof}

\paragraph{Reduction to a single binary relation.}
Here and below we use the following explicit signature encoding to justify the last sentence of the theorem. First replace each function, including a constant, by its graph relation under totality and uniqueness. Work on an inhabited sort $D$. A finite countermodel gives the actual graph tables. For the yes direction of the Horn reduction only a finite list of ground terms occurs in the constructed proof. Under the new graph axioms pick values successively for their subterms by totality. Use existential elimination to retain these fresh values locally and uniqueness to identify repeated occurrences. The translation of each original Horn axiom at those values gives precisely its ground consequence. More explicitly its antecedent universally quantifies over the named variable values. Its matrix says that for all consistent evaluations of the terms in that matrix the evaluated atoms satisfy the implication. For nullary facts use the same convention. The final atom $Y$ mentions none of the chosen values, so all the eliminations can be closed. This gives a PR intuitionistic proof compiler for the only direction used in Theorem~\ref{thm:church}.

For a finite relational signature use a directed edge $E$. Use parameter vertices $d$ (objects), $u_R$ (tuple nodes for relation $R$) and $p_{R,j}$ (coordinate nodes). A vertex $v$ has a stated type with tag $h$ if $E(v,h)$. Translate variables by relativizing quantifiers to $D(x):=E(x,d)$. Translate $R(x_1,\ldots,x_k)$ by
\[
 \exists v\left(E(v,u_R)\ \wedge \bigwedge_{1\le j\le k}\exists b\,
       (E(v,b)\wedge E(b,p_{R,j})\wedge E(b,x_j))\right).
\]
For $k=0$ the conjunction is empty. Equality is actual equality on $D$. Every graph and parameter tuple with inhabited $D$ now gives a decoded relational structure. Bounded scans give the truth identity by induction on formulas. For the converse encode any table by disjoint fresh vertices for all tags, all elements, every true tuple, and each of that tuple's coordinates. Draw just the specified type, incidence and value edges. Tag vertices have no outgoing edges. Each coordinate node links to its tag and its unique data element. Each tuple node links to its tag and its own coordinates. Among typed data the translation consequently recovers exactly the original tuples, including repetitions and nullaries.

For the separator use universal quantification over all tags, and an implication from inhabitance of $D$ and the relativized graph axioms and Horn assumptions to $Y^*$. A proof from an accepting trace runs verbatim inside the interpreted sort, using the inhabitance and totality hypotheses, then discharges and generalizes. In a rejecting graph the designated tags falsify that universal implication. For a finite-satisfiability reduction use existential tags, conjoin $\exists x\,D(x)$ and the translated sentence. The two decoding arguments give equivalence in both directions. Sorted inputs can first be relativized as in the finite-model chapter. All tags, formulas and finite incidence graphs are built by bounded loops.

\begin{corollary}
Intuitionistic and classical derivability are recursively enumerable and undecidable, even over the indicated binary signature. Finite validity is undecidable there as well. A decision procedure for ordinary classical validity, in its usual external reading, would likewise give the forbidden separator. There is no PR bound on the least proof code of a provable sentence as a function of its code.
\end{corollary}
\begin{proof}
The search in Section~\ref{sec:deduction} enumerates proofs. Finite soundness ensures that either proof calculus separates the yes and no cases if its theorem predicate were decidable. Finite validity separates them for the same reason. Ordinary semantics on all structures is sound for the displayed proofs and includes the finite countermodels. A PR upper bound, even one required only on provable inputs, would decide by bounded search for an accepted proof.
\end{proof}

\subsection{Recovering a real computation from a finite model}
\begin{theorem}[Finite satisfiability and Trakhtenbrot's obstruction]\label{thm:trakh}
One can compute a sentence $B_{P,w}$ in the language of a directed graph with equality with
\[
 \exists t\ \mathrm{Accept}(P,w,t)
 \quad\longleftrightarrow\quad
 \exists M\,[\mathrm{FinGraph}(M)\wedge V(M,B_{P,w})=1].
\]
There are PR witness maps both ways. Finite satisfiability is recursively enumerable, while finite validity has no recursive enumeration.
\end{theorem}
\begin{proof}
Describe the sentence first with two nonempty sorts, times and stack objects. Use a strict linear order on the times, a root $e$ of the stack sort, binary relations $U_a(v,v')$ for adding $a$ to stack $v$, two stack-valued functions $l(t),r(t)$, and a unary predicate $Q_i(t)$ for each label. Require the $Q_i$ to partition the times. Require $e$ to have no incoming $U$ edge and require every child to have at most one incoming parent-letter pair. Explicitly, $U_a(x,z)\wedge U_b(y,z)$ implies $x=y$ if $a=b$ and implies $\bot$ otherwise. The times $t_0,t_f$ are existentially specified least and greatest. The first label is start, $l(t_0)=e$, and $r(t_0)$ is reached from $e$ by a chain spelling the input in reverse order of its letters. Write this finite chain with existential intermediate objects. Require $Q_{\rm yes}(t_f)$.

Put $\mathrm{next}(t,t') := t<t'\wedge\neg\exists s(t<s\wedge s<t')$. For every adjacent pair require the disjunction of the legal instruction cases. For a push-left case its conjuncts are $Q_i(t),Q_j(t'),r(t')=r(t),U_a(l(t),l(t'))$. For pop-empty-left use the two corresponding labels, right-stack equality and $l(t)=e=l(t')$. For pop-$a$-left use its labels, right-stack equality and $U_a(l(t'),l(t))$. Include the symmetric right-stack cases. There are no disjuncts for transitions from a stopping label. These instructions expand to actual first-order formulas, of size bounded by a fixed polynomial in the written table and word if the two-sorted presentation is used.

An accepting run supplies time $0,\ldots,t$ up to its first stop, all its stack words and their suffixes, and the actual prefix edges between included words. Fill the two stack functions and each state bit from the trace. Every required conjunct follows from the successive configurations. In the converse direction sort the finite time table by its strict order: find the least unused element by bounded comparison, then mark it used. Every nonempty unused sublist has a least element, obtained by comparing its elements successively and retaining the smaller. Its successive choices exhaust the times and each consecutive pair satisfies $\mathrm{next}$. Carry, for each of the two current stack words, the list of object vertices from $e$ to the indicated top together with the letters. Initialization gives the first lists. A push appends the displayed edge. A pop from an empty word can only use the empty case because $e$ has no incoming edge. For a nonempty word the top has a recorded incoming edge, so differs from $e$, and uniqueness of the incoming parent and letter returns precisely the recorded predecessor. Delete that edge. Thus each move is precisely a machine step, and the last label certifies acceptance. Neither cycles outside the used paths nor unnamed stack objects affect the extraction. The induction stores words of length at most $|w|+$ the number of times and the two path lists. This is a bounded bit invariant of the time scan. Translate the functions, sorts and finite relational tables to a single graph by the preceding encoding. Reconstruction from graph to the time trace is a composition of PR operations.

Enumerating finite graph codes and evaluating a sentence semidecides finite satisfiability. For the final assertion suppose a machine enumerates exactly the finitely valid graph sentences, with an eventual output for each such sentence. Build $L$ that on input $u$ scans the enumeration for the sentence $\neg B_{u,u}$ and halts at yes just when it finds it. Set $u$ equal to the code of $L$ and write $H$ for acceptance of this run. By the construction and the assumed enumeration,
\[
 H\ \longleftrightarrow\ \neg B_{u,u}\text{ is finitely valid}
       \ \longleftrightarrow\ \neg H .
\]
The second equivalence uses finite Boolean evaluation and the two witness maps. From $H\to\neg H$ one gets $\neg H$, and then the other direction gives $H$, a contradiction. No case split on an unresolved halting question occurs.
\end{proof}

\section{Compactness, finite carriers and computable models}\label{sec:compactness}
The classical compactness theorem speaks of a single structure satisfying an entire theory whenever every finite subtheory has a model. Its finite-character direction is immediate from the definition of satisfaction. For derivations, the finite-support transformation gives a stronger conclusion: it returns the actual premises used. Replacing the model in compactness by a finite carrier, a computable presentation, or a certified member of a fixed list gives different arithmetical questions.

\begin{proposition}[Two separations for finite satisfiability]
There is a single sentence true in the natural-number successor structure and false in every finite structure. There is also a primitive-recursive sequence of sentences for which every finite sublist has a uniformly constructible finite model, while every particular finite structure falsifies a computable member of the sequence.
\end{proposition}
\begin{proof}
In the language of $c,s$ take
\[
 (\forall x\forall y\,(s(x)=s(y)\to x=y))\ \land\
 (\forall x\,s(x)\ne c).
\]
Its direct interpretation $c=0$, $s=S$ is a conjunction of axioms of iPRA. In a finite model of the first conjunct, the $m$ table entries of $s$ on $[m]$ are distinct. Mark each output in an $m$-bit array. After $j$ inputs the array has exactly $j$ marked entries by induction on $j$, since the next output has not yet occurred. All $m$ entries are therefore marked, including $c$, which contradicts the second conjunct. This fixed interpretation requires no global truth predicate for arithmetic.

For the sequence take the pure-equality sentences $H_j$ asserting the existence of $j$ pairwise distinct elements, with $H_0=\top$. Their codes are obtained by two bounded loops for the matrix and one for the prefix. For a finite sublist use the maximum of $1$ and its indices as the carrier size. Use the first $j$ elements as witnesses for $H_j$. A carrier of size $m$ fails $H_{m+1}$ by the same marking argument. Thus even a supplied model for every finite stage does not ensure a common finite model.
\end{proof}

\begin{theorem}[Exact finite obstruction size at a fixed cardinal bound]\label{thm:boundedcover}
Fix a nonempty size range $1,\ldots,B$ and a finite signature. Let $K$ be the number of isomorphism types in that range. A finite list of sentences having no joint model of size at most $B$ has a computable sublist of size at most $K$ with the same property. The constant $K$ is sharp and is PR in the signature and $B$.
\end{theorem}
\begin{proof}
Enumerate all labelled structures, numbering a function row by a base-$m$ choice and a relation row by a bit. The list length is
\[
 \sum_{m=1}^B
  \left(\prod_f m^{m^{\operatorname{arity}(f)}}\right)
  \left(\prod_R 2^{m^{\operatorname{arity}(R)}}\right).
\]
Keep a structure if it is the first in its isomorphism class. To test isomorphism between two size-$m$ tables, scan the $m^m$ candidate maps, testing bijectivity and every operation and predicate row. Distinct sizes need no test. This gives $K$ representatives with an explicit PR loop bound. In each representative find the first sentence of the input list it falsifies. The assumed failure of joint satisfaction ensures success of each search. Retain these sentences, deleting duplicates. Each candidate structure is isomorphic to its representative, so Theorem~\ref{thm:maps} transfers the falsity of the selected sentence. This proves the upper bound even when the original family contains repeated sentences.

For sharpness form one sentence $\neg D_M$ from Theorem~\ref{thm:diagrams} for each representative $M$. The whole list eliminates every type. If one member $\neg D_M$ is removed, $M$ satisfies every remaining member. Thus no smaller universal bound works. For an enumerated theory a supplied list of failing axiom indices, one per representative, similarly gives a single inconsistent-with-the-bound prefix by taking their maximum. This last statement concerns a supplied finite certificate. It makes no infinite selection of failure indices.
\end{proof}

\begin{theorem}[Computable finite approximations with no computable common model]\label{thm:nocompmodel}
In the fixed language of a constant $c$, unary function $s$, and unary predicate $P$, one can give an increasing PR sequence of finite theories $T_k$ and a PR finite model $M_k$ for each $T_k$. No total computable bit function can supply the joint values of all the ground atoms $P(s^j(c))$ compatibly with all these theories. In particular their union has no model with a computable atomic diagram and computable interpretations of $s,c$. Its finite-proof consistency is provable in iPRA.
\end{theorem}
\begin{proof}
Use an ordinary deterministic program enumeration with PR bounded simulation. Let $U(e,j,t)$ record whether program $e$ has halted on input $j$ within $t$ steps, and its output when it has. Halting is made a stationary state. The scan up to $t$ is primitive recursion on configurations. In particular a reported output persists for every greater bound. Write $A_j=P(s^j(c))$. For each $e<k$, inspect $U(e,e,k)$. On output bit $b$ include in $T_k$ the literal forcing $A_e$ to have bit $1-b$. Thus on output $0$ include $A_e$, on output $1$ include $\neg A_e$, and on any other outcome include neither. Treat finite theories as lists up to inclusion of members. Stationarity gives $T_k\subseteq T_{k+1}$.

For $M_k$ use $[k+1]$, with $c=0$ and $s(i)=\min(i+1,k)$. Set $P(i)$ to the forced bit if the bounded simulation produced one for $i<k$, and to zero otherwise. For $i<k$ the value of $s^i(c)$ is $i$, by the induction identity $s^j(c)=\min(j,k)$. Hence $M_k$ satisfies every listed literal. Any finite derivation from the union, with its theory-premise occurrences certified by stages, has a maximum cited stage $k$ (take zero for the empty list). The model $M_k$ validates all its premises. Finite two-valued soundness precludes a conclusion of falsity. This proves the numerical statement that no certified contradiction is accepted.

Suppose program $e$ purports to compute the sequence of bits of a common model on the $A_j$. If it is total with bit outputs, applying its totality hypothesis just at $j=e$ gives a halting time $t$ and bit $b$. Set $k=\max(e+1,t)$. The stage-$k$ axiom for $A_e$ forces $1-b$, while the proposed computation gives $b$. In particular this refutation takes an explicit halting trace as input, with a PR bound on the stage needed. Computable $c,s,P$ would compute the forbidden sequence by iterating $s$ exactly $j$ times before querying $P$. Externally, in classical set theory, an ordinary model can be described on $\mathbb N$ by taking $c=0$, $s=S$ and declaring $P(e)$ true just when program $e$ on input $e$ eventually outputs $0$. The arithmetical obstruction concerns computable truth bits for a common model, as requested by an effective strengthening of compactness.
\end{proof}

\begin{proposition}[A branch encoded by a common atomic diagram]\label{prop:compbranch}
For every fixed PR prefix-closed binary tree whose every level is inhabited, a uniform construction gives a sequence of finitely satisfiable stages in the same three-symbol language such that the ground truth bits of any common model form an infinite path. Conversely any supplied path gives an interpretation on the successor chain satisfying all the stages.
\end{proposition}
\begin{proof}
A binary word $w$ of length $h$ is forbidden if its tree test is zero. At stage $k$ enumerate all forbidden words of lengths $h\le k$. For each include the clause saying that some $A_j$, $j<h$, differs from $w_j$. Use the positive literal $A_j$ if $w_j=0$ and the negative one otherwise. For the empty word the clause is falsity. At level $k$ a bounded scan of the $2^k$ words finds the first tree node $v$ by the hypothesis. Its prefixes belong to the tree, so $v$ differs from each forbidden word at some coordinate in that word. Interpret $A_0,\ldots,A_{k-1}$ as the bits of $v$ in the truncated successor chain $[k+1]$ just used. All stage clauses hold. If the bits of a common model had a forbidden prefix $w$, its stage-$h$ clause would be false, term by term. Conversely a supplied path satisfies every clause for this same reason on $\mathbb N$ with $P(j)$ its $j$th bit. For a PR path this is a schema in the supplied bit function. For an arbitrary path the converse describes the external interpretation. Finite verification uses at most $\sum_{h=0}^{k}2^h$ word tests per stage plus bounded coordinate comparisons. This reduction identifies the infinite branch that a general model-construction principle would have to provide.
\end{proof}

\begin{theorem}[A Herbrand tree with an exact finitary completeness statement]\label{thm:herbrandtree}
Let $S$ be a finite list of universally closed clauses over a finite equality-free signature having at least one constant. There is a PR binary tree $\mathcal T_S$ for which iPRA proves
\[
 \neg\exists p\,\operatorname{GroundRef}(S,p)
 \quad\longleftrightarrow\quad
 \forall k\,\operatorname{HasLevel}(\mathcal T_S,k)=1 .
\]
Here a refutation is a ground resolution tree with the source clause and substituting ground terms attached to each leaf. Both a missing level's refutation and a refutation's missing level are primitive-recursively recoverable. A supplied infinite path yields a Herbrand model, that is, a model whose elements are the ground terms and whose functions build those terms.
\end{theorem}
\begin{proof}
Use natural-number syntax codes themselves as coordinates. A coordinate which is not the code of a ground atom has no semantic role. This convention works even when there are finitely many ground atoms or none. For a word $w$ of length $k$, inspect every $i<k$. If $i$ codes a pair consisting of a valid index into $S$ and an appropriate tuple of well-formed ground terms, form the indicated ground instance. If every atom code of the instance is less than $k$, require at least one of its signed literals to be true under $w$. Otherwise impose no condition from $i$ at this stage. In particular an inspected empty clause causes failure. Declare $w\in\mathcal T_S$ when all these bounded requirements succeed. Every requirement tested at a shorter length was already tested at the larger length and reads the same coordinates. Thus prefixes of tree members are members. The level test scans all $2^k$ words, recording the first success if there is one.

If level $k$ is missing, collect the inspected ground instances whose atom codes are below $k$. No valuation satisfies this finite family: every assignment to its atoms extends to a $k$-bit word by assigning zero to unused positions. The bounded resolution construction in Section~\ref{sec:normal} returns a refutation with at most $2^{k+1}-1$ nodes, and each selected leaf retains its generating code $i$. Conversely, given a refutation take $k$ greater than all its leaf generating codes and all atom codes used anywhere in it. If a length-$k$ word passed the test, it would satisfy all the leaves. Resolution soundness would make its empty root clause true. Therefore this level is missing. These two PR maps prove the displayed equivalence by the decidability of the level test. In particular ordinary syntactic consistency of $S$ implies every level is inhabited: universal instantiation proves the leaves from $S$, and the resolution rule on disjunctions follows by cases on the two premises, using contradiction in the two opposite-pivot cases taken together.

For a supplied path $b$, let the domain consist of the well-formed ground term codes. A constant ensures it is inhabited. Interpret $f$ by the constructor on its term arguments, and $R(\bar t)$ by the bit $b$ at the code of the corresponding ground atom. Evaluating a term under a tuple of ground terms gives its syntactic instance, by term induction. Given a clause and an assignment, compute the generating index $i$ of that instance and choose a level exceeding $i$ and every atom code in the instance. The corresponding path prefix certifies its truth. This proves every universal clause. Conversely, in an ordinary model with Boolean ground truth values, evaluate each ground atom to supply a path bit and use zero on unused coordinates. Each inspected clause holds by universal instantiation. For a supplied PR path the forward construction is a schema about PR atomic tables on coded terms. The theorem never needs an operation choosing a common path from its inhabited levels.

If there is logical equality, first use a fresh binary predicate for equality and add universal clauses for reflexivity, symmetry, transitivity and coordinatewise congruence for each original symbol (as implications of conjunctions of atoms, written as clauses). The tree construction applies verbatim to this expanded list. Along a supplied path every term $t$ has a canonical representative, the least ground term code $u\le t$ declared equal to $t$. Reflexivity makes the search succeed. Equality of representatives agrees with the path's equivalence relation by symmetry and transitivity. Applying a constructor and replacing its result by this representative, and reading predicates on representatives, now gives the genuine-equality quotient interpretation. The proof is the projection argument of Theorem~\ref{thm:eqquot}. Here each particular term's representative search is bounded by that term's own code rather than by a finite carrier size.
\end{proof}

\section{Lean formalization and the remaining object-theory proofs}\label{sec:lean}
The accompanying Lean files are standalone developments using Lean 4.32.0 and \texttt{Std}. Their theorems have the following scopes.
\begin{description}[style=nextline,leftmargin=1.25em]
\item[\texttt{NegativeTranslation.lean}]
\texttt{FOLNegative.translate} and \texttt{recover} construct typed derivations in both directions, with de Bruijn capture avoidance, the ordinary intuitionistic connectives and quantifiers, equality transport, and variadic function and relation symbols. \texttt{negative\_translation\_iff} and \texttt{equiconsistency} state the resulting equivalences. The raw syntax admits any numerical index or argument list. For use at a fixed signature the paper's range and arity tests restrict it to well-scoped input. The inductive calculus makes admissible context weakening and individual shifting explicit. \texttt{finite\_sound} covers finite inhabited carriers for both calculi. \texttt{consistency} excludes an empty-context deduction of falsity.
\item[\texttt{Substitution.lean}]
\texttt{Formula.truthSubst} proves semantic substitution for an executable Boolean interpretation with finite carriers, equality, arbitrary relation argument lists, and function symbols of arities zero, one and two. Scope is intrinsic via \texttt{Fin}. Renaming, coincidence on free indices and identity substitution are proved with binder lifting. Extending this particular file to all function arities adds mutual recursion on argument lists.
\item[\texttt{BackAndForth.lean}]
\texttt{FirstOrderPacket.transfer} proves rank invariance for arbitrary related environments, and \texttt{sentence\_transfer} specializes to sentences. It uses scoped relational atoms with variable arguments, equality, falsity, conjunction, negation and existence. Its truth clauses use Lean propositions. Duplicator's reply clauses transport existential witnesses in both directions. The Python companion implements the converse separator on finite tables.
\item[\texttt{GroundResolution.lean}]
\texttt{check\_sound} proves that every accepted resolution tree preserves a valuation. \texttt{refute\_checked} proves that a finite family of clauses having no Boolean model has a constructed tree accepted with the empty root. The implementation builds an assignment tree on a bounded set of atoms, then drops unused pivots. The checker accepts clause lists with duplicates.
\item[\texttt{CanonicalForcing.lean}]
\texttt{canonical\_completeness} proves the propositional part of Section~\ref{sec:canonical}, including falsity at inconsistent syntactic worlds. Worlds are finite contexts and implication ranges over larger contexts. The first-order extension uses substitution arrows and the generic individual in the human proof.
\end{description}
The files present different interfaces for raw proof trees, intrinsically scoped expressions and finite tables. A common numerical decoder is needed to relate their Lean types.

\subsection{An object-calculus example}
As a separate test of object-level proof production, the archive includes a native scalar certificate for
\[
     x=y,\quad y=z\quad\vdash\quad H(x)=H(z),\qquad
     H(a)=I(a,S(a);u.I(u,a;v.S(v))).
\]
Here $I$ denotes ProveIt's scalar iterator, with the variable following the semicolon bound in the step body. The proof combines equality transport with the iterator computation axiom. Both free occurrences and occurrences under an iterator have to be treated correctly. ProveIt's pinned raw goal checker accepted this certificate. The certificate, checker revision and reproduction procedure are in the \texttt{native} directory. This is one checked native instance of the equality calculus.

\subsection{The object language and numerical interface}
Use the upstream namespace \texttt{LeanProofs.ConciseIPRA}. Its \texttt{PR k} constructors are zero, successor, projection, composition, and primitive recursion (recursion argument first). \texttt{PTm n} applies these symbols to terms with $n$ free slots. \texttt{PRAAx} supplies the defining equations, the successor axioms and quantifier-free induction. A completed internal theorem is a term of type \texttt{Prf PRAAx []} applied to its explicit universally closed arithmetic formula. A Lean theorem on \texttt{Nat} or an inductive derivation type does not automatically supply this term.

Use tagged lists for records with the pairing of Lemma~\ref{lem:toolkit} or a sentinel variant of the repository pairing. For the latter store cons as $1+\langle a,b\rangle$ with $\langle a,b\rangle=(a+b)^2+a$, reserve zero for nil, and reject a positive value unless its projections re-encode it. The pairing has gaps. Components of a cons are strictly smaller than its code. Use the following layouts, rejecting malformed fields.
\begin{center}\small
\begin{tabularx}{\linewidth}{@{}lX@{}}\toprule
Record & Fields and tests \\ \midrule
Signature & Lists of function and relation arities, in index order.\\
Syntax node & Constructor tag, symbol or variable index as needed, and ordered children. Expected sorts and scoping follow Section~\ref{sec:binding}.\\
Proof node & Mode, scope, context, conclusion, rule, annotation and children. Each rule selects its term, equality body or assumption position and checks the required premise sequents.\\
Structure & Signature, positive cardinality and tables with exact lengths and value ranges. An assignment is a list of domain elements of stated length.\\
Kripke structure & World and label counts, preorder bits, inhabited domain masks, persistent atom and coherent function tables on active tuples.\\
EF position & Two relational structures, two equal-length assignments and remaining rounds. A separator is another scoped formula.\\
Computation & Program table, input word, and list of triples of label and two stacks. Test the first record and consecutive steps, then read the stopping label.\\
\bottomrule\end{tabularx}\end{center}
The decoder outputs a typed tree or failure. Prove that decoding the encoding of a scoped tree gives that tree, and that re-encoding a canonical valid input gives the same number. Both use a postorder scan and pairing equations. Fix unused fields to a single default, or normalize them before the second test.

\subsection{Compiling the inductions}
Associate to each fixed operation $f$ a symbol $\widehat f$ of the correct \texttt{PR} arity. Projections and tuple composition implement parameter passing. A bounded loop with initial state $a(\bar x)$ and step $h(i,u,\bar x)$ compiles to \texttt{PR.prec} on their codes. Use a tuple state for multiple outputs. The repository's pairing, simultaneous recursion and finite-parameter search lemmas provide checked components. Its \texttt{Search.Finite.mu} searches for a \emph{zero} below the bound and returns the bound on failure. To search for a one of a bit $h$, use $1\mathbin{\dot-}h$.

For example a semantic row stores $h(i)=[B](i::e)$. Set $j=\operatorname{mu}(1\mathbin{\dot-}h,m,e)$ and define the existential bit by $[j<m]$. From $i<m$ with $h(i)=1$, \texttt{search\_found} gives $j<m$ and a zero of the complemented body. Conversely, \texttt{search\_spec} at $j$ supplies that zero from $j<m$, and its earlier-index condition proves minimality. The bit equations for $h$ recover $h(j)=1$. This supplies the derivation of the semantic existential clause and a witness uniform in the body row.

For a loop with parameters $x$, introduce $u_i$ for its PR state and write the induction formula as an equation $\widehat{\mathrm{ok}}(i,u_i,x)=1$. At an EF level this bit ANDs all completed position checks, including a losing move with a correctly oriented separator or a winning position with the required replies. For rank preservation on an input formula, form a second finite table over its subformulas, positions and remaining budgets. The input formula is a free parameter of this induction. The collection of all bounded-rank formulas is infinite. With the initialization and update equalities derived, apply \texttt{PRA.qfInd} and specialize at the final index. Universally close the original input variables last.

Representative target formulas follow. Close the input variables universally and guard the equations by valid signatures, models, contexts and positions of the indicated sorts.
\begin{align*}
 \mathrm{Check}_C(p,\Gamma,A)=1
   &\ \to\ \mathrm{Check}_I(\mathrm{translate}(p),\Gamma^N,A^N)=1,\\
 \mathrm{Well}(M,e,A,\sigma)=1
   &\ \to\ V(M,e,A[\sigma])=V(M,e_\sigma,A),\\
 W_r(u,v)=0
   &\ \to\ \mathrm{Distinguishes}_{\le r}(\mathrm{sep}(r,u,v),u,v)=1,\\
 \mathrm{Trace}_{\rm yes}(P,w,t)=1
   &\ \to\ \mathrm{FinModel}(\mathrm{encode}(P,w,t),B_{P,w})=1,\\
 \mathrm{FinModel}(M,B_{P,w})=1
   &\ \to\ \mathrm{Trace}_{\rm yes}(P,w,\mathrm{extract}(P,w,M))=1.
\end{align*}
A history and its stopping bound are interconvertible in these last two lines by bounded simulation and length. The distinguishing test conjoins scope, rank and the two evaluated bits. Each endpoint above is an implication between quantifier-free equalities.

\paragraph{Remaining dependencies.}
After pairing and row access, prove shift/composition and the fold laws, then the two substitution lemmas and the local rule checks. Use them for grafting and forcing. Build Skolem rows and resolution after finite evaluation. Use resolution and congruence representatives for Bernays--Schönfinkel and the Herbrand tree. For witness closure, verify the finite mask fixed-point count and the representative round trip for closed terms. Use EF replies for monadic truncation, order thresholds and cycles, and the symbolic elimination templates for monadic refutations. For computation reductions verify two-stack steps, Horn instantiation, rejecting tables, sorted time extraction and graph decoding in that order. Represent machine totality by $\forall x\exists t\,\mathrm{StopsBit}(e,x,t)$ and specialize it at the diagonal input. For canonical forcing, a Lean recursion on the fixed template constructs a separate \texttt{Prf} goal and derivation, instantiating the child derivations at each translated arrow. The uniform embedded-iPRA proofs of the full suite are remaining formalization work.

\clearpage
\section{Source guide, boundaries and further questions}
\begingroup\small
The starting article is \href{https://en.wikipedia.org/wiki/First-order_logic}{\emph{First-order logic}}. The table locates the mathematical claims selected from its sections and linked topics. The additional bounds, explicit algorithms and proof maps are developed in this manuscript.

\begin{center}\small
\begin{tabularx}{\linewidth}{@{}>{\raggedright\arraybackslash}p{.32\linewidth}>{\raggedright\arraybackslash}X@{}}\toprule
Source topic & Result obtained here \\ \midrule
Formation, variables and substitution & Unique readability, alpha-equivalence, capture avoidance, composition and semantic substitution (Sections~\ref{sec:binding}, \ref{sec:models}).\\
Deductive systems, equality, \href{https://en.wikipedia.org/wiki/Extension_by_definitions}{extension by definitions} & PR proof checking, deduction, finite support, equality transport, quotient structures and conservativity.\\
Intuitionistic first-order logic & Increasing-domain soundness and countermodels, embedding classical derivations, and canonical forcing for the negative fragment.\\
Interpretations, variations and \href{https://en.wikipedia.org/wiki/Relational_algebra}{relational algebra} & Finite satisfaction, preservation and diagrams, elimination of functions and sorts, relational operations and fragment-elementary substructures.\\
\href{https://en.wikipedia.org/wiki/Prenex_normal_form}{Prenex} and \href{https://en.wikipedia.org/wiki/Skolem_normal_form}{Skolem} normal forms, automated deduction & Witness tables, two clausal forms, resolution completeness and two-way cut-free Herbrand extraction.\\
Decidable fragments & Monadic rank profiles and proof-producing decision with equality, and the relational existential--universal fragment with constants and equality.\\
\href{https://en.wikipedia.org/wiki/Ehrenfeucht-Fraisse_game}{Ehrenfeucht--Fraïssé games} and expressive limitations & Formula separators, exact finite-order threshold, parity, connectedness and reachability counterexamples.\\
Completeness and undecidability & Church separation for one binary relation, finite satisfiability, nonenumerability of finite validity, and an arithmetical completeness reading.\\
\href{https://en.wikipedia.org/wiki/Compactness_theorem}{Compactness} and \href{https://en.wikipedia.org/wiki/Lowenheim-Skolem_theorem}{Löwenheim--Skolem} & Sharp bounded-carrier obstruction, no computable common diagram, Herbrand trees, and Skolem hulls with supplied witness selectors.\\
\bottomrule\end{tabularx}\end{center}

\paragraph{ProveIt.}
The \href{https://github.com/VladimirReshetnikov/ProveIt/tree/bbb85238dac533fc4b7b84a750cca864fe35fed5/Logic/PrimitiveRecursiveArithmetic}{iPRA project} fixes the base theory. Its ordinary deduction target is specified by \href{https://github.com/VladimirReshetnikov/ProveIt/blob/bbb85238dac533fc4b7b84a750cca864fe35fed5/Logic/PrimitiveRecursiveArithmetic/ConciseIPRA/Lean/ConciseIPRA/PR.lean}{\texttt{ConciseIPRA/PR.lean}}. Arithmetic providers include \href{https://github.com/VladimirReshetnikov/ProveIt/blob/bbb85238dac533fc4b7b84a750cca864fe35fed5/Logic/PrimitiveRecursiveArithmetic/Derivations/Lean/IPRADerivations/Core/Pairing.lean}{pairing and simultaneous recursion} and \href{https://github.com/VladimirReshetnikov/ProveIt/blob/bbb85238dac533fc4b7b84a750cca864fe35fed5/Logic/PrimitiveRecursiveArithmetic/Derivations/Lean/IPRADerivations/Core/SearchParameters.lean}{bounded search with parameters}. The separate \href{https://github.com/VladimirReshetnikov/ProveIt/tree/bbb85238dac533fc4b7b84a750cca864fe35fed5/Logic/FirstOrder}{first-order completeness project} reports an independent Lean and Rocq development for equality and one binary relation. Its Lean model-existence proof uses classical choice. Its general-signature semantic compactness target is a separate mathlib-backed result.

\paragraph{Comparison literature.}
Y. Forster, D. Kirst and D. Wehr, \href{https://arxiv.org/abs/2006.04399}{\emph{Completeness Theorems for First-Order Logic Analysed in Constructive Type Theory}}, compare completeness principles for different semantics and study connections to weak König's lemma. D. Kirst and D. Larchey-Wendling, \href{https://arxiv.org/abs/2104.14445}{\emph{Trakhtenbrot's Theorem in Coq: Finite Model Theory through the Constructive Lens}}, give a signature-sensitive classification of finite satisfiability. Here both ends of the two-stack reductions are explicit arithmetic codes. The proof transformation in Section~\ref{sec:negative} is the Gödel--Gentzen construction. See J. Avigad and S. Feferman, \href{https://math.stanford.edu/~feferman/papers/dialectica.pdf}{\emph{Gödel's functional (Dialectica) interpretation}}, for background.

\paragraph{Scope of the reformulations.}
Arbitrary uncountable signatures, the assertion that every consistent theory has a set-sized Tarski model, the upward cardinal conclusion of Löwenheim--Skolem, and the characterization of abstract logics in Lindström's theorem involve objects outside the numerical language chosen here. The computable-diagram counterexample explains a precise obstruction to replacing models by programs. Other arithmetizations of completeness require their own analysis. The extraction in Section~\ref{sec:normal} accepts cut-free input. The guarded and two-variable fragments and the full classification of decidable signatures are further topics beyond the decisions proved here.

\paragraph{Further research questions.}
Can the binary-signature separation be formalized in ProveIt with a reusable theory of interpreting finite signatures in graphs? Can the monadic decision proof exploit shared subproofs to give a substantially smaller verified bound on certificate size? Which additional fragments admit constructive completeness with a PR certificate transformation for their intended semantics?

\endgroup

\end{document}
